\chapter{体論}
この章では体論・ガロア理論についてまとめます. 代数学IIでは, 分離拡大や正規拡大など, 体拡大に関する基本事項を学び, とくにガロア拡大に現れる強い対称性を調べました. また, 有限次ガロア拡大の中間体とガロア群の部分群を対応させる, ガロア理論の基本定理を学習しました. 体は一般に無限集合であり複雑に見えますが, ガロア群の位数は拡大次数と一致するため, 有限群論の枠組みで調べられます. このため, 有限群論的な現象と体拡大に関する現象を相互に対応させることができます. \par
体拡大としては$\mb{Q}$上の拡大(代数体)に関することが多いが, もっと特殊な体拡大(たとえば有限体, 円分体など)に関することも学習した. 円分体のガロア群は$(\mb{Z}/n\mb{Z})^{\ast}$の形であり, 有限体の拡大はフロベニウス写像によって生成される巡回群である, という事実もたいへん重要である.
\par
問題演習においては, $\mb{Q}$上の多項式の最小分解体を調べ, そしてそのガロア群について調べるという問題が非常に多かった. まず拡大次数を求めることが鉄則だが, そのためには$\mb{Q}$上の多項式の既約性を求めることがよくある. そのためにEisensteinの既約判定法やガウスの補題など, 環論的ではあるがさまざまな命題について学習する必要がある.



\section{基礎事項}
\subsection{既約多項式・最小多項式}

\tf{参考}
\ben
\item 雪江整数論1, 8章2節, 241p~
\een


体拡大の次数を求めるには, 多項式の既約性を調べることがしばしば必要である.

次は高校数学でも知られている可約判定である.
\begin{prop}
$f(x) = \sum_{i=0}^{n} a_i x^{n}\in \mb{Z}[x]$ ($a_n\neq 0$)とする.このとき, $f(x)$の有理数の根はかならず
\[\pm \dfrac{(\text{$a_0$の約数})}{\text{$a_n$の約数}}\]
の形をしている.とくに$f$がモニック($a_n=1$)ならば, 有理数根は必ず整数である.
\end{prop}
\lefba{
\begin{prop} (\tf{Eisensteinの既約判定法})
$A$をUFD(たとえば$\mb{Z}, \mb{C}[X,Y]$など), $p$を素元として.  $f(T)\in A[T]$が次の条件を満たすとする:
\[f(T) = T^n + a_{n-1}T^{n-1} + \dots + a_1 T + a_0\]
と表したとき, $p|a_i$ ($0\leq i\leq n-1$)かつ$p^2\not | a_0$. \par
このとき, $f(T)$は$A[T]$において既約である.
\end{prop}
}
上は便利だが, 万能ではない.次の方法があることも覚えておくと良い.
\begin{prop}
$A$を正規環, $K$を$A$の商体, $\maf{p}\subset A$を素イデアル, $f(x) \in A[x]$をモニック, $\deg{f(x)} = n$とする.このとき, $f(x)$の係数を$\maf{p}$を法として考えた多項式$\ovl{f}(x)$が$(A/\maf{p})[x]$で$m<n$次の因子を持たなければ, $f(x)$も$m$次の因子を持たない.とくに$\ovl{f}(x)$が$A/\maf{p}$の商体上既約なら$f(x)$も既約である.
\end{prop}
\begin{cor}
$p$を素数, $f(x)\in \mb{Z}[x]$をモニック, $\deg{f(x)} = n$とする.$f(x)$の係数を$\bmod{p}$した多項式$\ovl{f}(x)$が$\mb{F}_{p}[x]$で$m<n$次の因子を持たなければ, $f(x)$も$m$次の因子を持たない.とくに$\ovl{f}(x)$が$\mb{F}_{p}$上既約なら$f(x)$も既約である.
\end{cor}
次も重要な補題である.$\mb{Z}$の範囲のみで済むというのが非常に効果を発揮する場面がよくある.
\tbox{
\begin{prop} $A$を\tf{UFD}, $f,g\in A[x]$を原始多項式, $K$を$A$の商体とする.
\ben
\item $fg$も原始多項式である.
\item $f$が$A[x]$で既約$\iff$ $f$が$K[x]$で既約.
\een
\end{prop}
}{title=ガウスの補題}
\begin{eg} ガウスの補題により, 既約性が保証される範囲を広げることができる.
\ben
\item $A=\mb{Z}$のとき$K=\mb{Q}$である. $f(x) = x^2 - 2\in \mb{Z}[x]$は原始多項式である.  $f(x)$が$\mb{Z}[x]$で既約であることを示す. 可約なら $f(x) = (ax+b)(cx+d)$と書けるはずだが, $ac = 1$より$a=c=\pm 1$で, $a=c=1$としてもよい. このとき$b+d=0$で, $bd = -b^2 = -2$となるが, このような$b$は無いので矛盾. したがって$f(x)$は$\mb{Z}[x]$において既約であり, ガウスの補題から自動的に$\mb{Q}[x]$でも既約であることが従う.
\item $A=\mb{C}[X]$のとき$K=\mb{C}(X)$である. $\mb{C}[X,Y]$を$\mb{C}[X][Y]$と見る. $f(Y) = Y^2 - X^3$は$A[Y]$において既約である(上と同様にせよ). よって, $f(Y)$は$\mb{C}(X)[Y]$において既約である. すなわち, いかなる有理関数$p(X),q(X),r(X),s(X)$を用いても$(p(X)Y+q(X))(r(X)Y + s(X))$と表すことは不可能である.
\een
\end{eg}

\tbox{
\begin{prop}
$A$をUFD, $\mr{Frac}A = K$とする. $\alpha \in K$の$K$上最小多項式を
\[F(T) = T^m + k_1T^{m-1} + \dots + k_{m-1}T + k_m\]
とする($k_i\in K$). さらに$\alpha$が$A$上整で, その整等式の中で最小次数のものを
\[G(T) = T^n + a_1T^{n-1} + \dots + a_{n-1}T + a_n\]
とする($a_i\in A$). このとき, \bolm{F(T) = G(T)}である.
\end{prop}
}{title=整元の最小多項式は環のほうの係数で取れる}
\prf $F$は$K$上の最小多項式であるから, $K[T]$において $F(T)|G(T)$である. $G(T) = F(T)H(T)$と表す($H(T) \in K[T]$). $F,H$の係数の最小公倍数を払って原始多項式$\witi{F}(T), \witi{H}(T)$を得るとする. $a,b\in A$に対して, $\witi{F}(T) = aF(T)$, $\witi{H}(T) = bH(T)$とおく. このとき$ab\cdot G(T) = \witi{F}(T)\witi{H}(T)$で, 右辺は原始多項式なので左辺は原始多項式. $G(T)$はモニックなので原始的であり, $ab\in A^{\times}$である. よって$a,b\in A^{\times}$である. \ao{よって, もともと$F,H$の分母は無かった. }したがって$F,H\in A[T]$であり, とくに$F$は$A$係数. $G\in A[T]$の最小性より$G(T) | F(T)$であり, ともにモニックなので $F(T) = G(T)$が従う

\qed





\begin{prac}
$k$を体とする. $k[X,Y]/(Y^2 - X^3)$が整域であることを示せ. (\tf{Hint}: $Y^2 - X^3$が"既約元"であることを示せば, 素元となるので素イデアルによる剰余環と見れる. )
\end{prac}


\begin{prac} (京大の\tf{学部入試})\\
有理数係数多項式$P(x)$が$P(\sqrt[3]{2})=0$を満たすならば, $P(x)$は$x^3-2$で割り切れることを証明せよ.
\end{prac}

\begin{egprob}
次の多項式が$\mb{Q}$上既約であることを示せ.
\ben
\item $f(x) = x^{p-1} + x^{p-2} + \cdots + x + 1$ ($p$は素数)
\item $f(x) = x^4 - x - 1$
\item $f(x) = x^4 + x + 1$
\een
\end{egprob}
\begin{egprob}
$x=t+1$とおく.$(x-1)f(x) = x^{p} - 1$から$tf(x) = (t+1)^{p} - 1$である.右辺は$t(t^{p-1} + pt^{p-2} + \cdots + {}_{p}\mr{C}_{2} t^2 + p)$だから$f(x) = t^{p-1} + p t^{p-2} + \cdots + p$でEisensteinの既約判定法からこれは$t$の多項式とみて既約である.よって$x$の多項式とみても既約である.\qed
\end{egprob}

\begin{prob}
次の多項式が$\mb{Q}$上既約であることを示せ.
\ben
\item $f(x) = x^3 - 105x^2 + 147x + 63 $
\item $f(x) = x^4 + x^2 - x  + 2$
\item $f(x) = x^5 - x - 1$
\een
\end{prob}
\begin{prob}

\end{prob}


根の最小多項式を求める場面もあります. このとき, 試行錯誤によって多項式を推測せずに求める方法があります.
\begin{egprob} (\href{https://www.math.kyoto-u.ac.jp/~yukie/Alg-II-final.pdf}{代数学II期末試験(年度不詳)})
$f(x) = x^3 - 2x + 2$の一つの根を$\alpha$とする. $\beta:= \alpha^2 + \alpha$の最小多項式を求めよ.
\end{egprob}
\begin{egans} $K=\mb{Q}(\alpha)$とする. $f$はEis多項式なので既約. よって$[K:\mb{Q}] = 3$. $\set{1,\alpha,\alpha^2}$は$K$の$\mb{Q}$上の基底. この基底に関し, $\beta$倍写像の行列表示を求める.
\bealn{
\beta \alpha &= \alpha^3 + \alpha^2 = \alpha^2 + 2\alpha - 2 \\
\beta \alpha^2 &= 2\alpha^2 -2
}
より, $\beta$倍写像は
\[
\bmat{
0&-2&-2 \\
1&2&0\\
1&1&2
}
\]
であり, これの特性多項式が$\beta$を根に持つはずである. 計算するとそては$x^3 - 4x^2 + 8x - 6$になり, Eis多項式である. よってこれが答え.
\end{egans}

たまに次も使える.


\begin{prop}
$f(x)$を$K$上の分離多項式とする. $f(x)$の根を$\alpha_1,\dots, \alpha_n$とする. $L=K(\alpha_1,\dots, \alpha_n)$とする.
\ben
\item $f(x)$は$K$上既約.
\item $\gal{(L/K)}$は集合$\set{\alpha_1,\dots,\alpha_n}$に推移的に作用する.
\een
\end{prop}




\subsection{拡大次数}
$L/K$が体の拡大であるとする.包含写像$K\to L$によって$L$は$K$代数(環の構造を持った$K$ベクトル空間)と考えられる.このベクトル空間としての次元を$L/K$の\tf{拡大次数}といい, $[L:K]$で表す.
\begin{prop} (\tf{拡大次数まとめ})
\ben
\item \bolm{[L:K] = [L:M][M:K]}
\een
\end{prop}

\begin{prac}
$n,m$を自然数とする.$\mb{Q}(\sqrt{m})\cap \mb{Q}(\sqrt[3]{n}) = \mb{Q}$を示せ.($m,n$で場合分けする必要が少しある)
\end{prac}
\begin{prac}
$[\ovl{\mb{Q}}:\mb{Q}] = \infty$を示せ.
\end{prac}

\begin{eg}
$p$を素数とする.$\mb{Q}(\cos{\dfrac{2\pi}{p}})/\mb{Q}$と$\mb{Q}(\sin{\dfrac{2\pi}{p}})/\mb{Q}$の拡大次数を求める.
\end{eg}

\subsection{代数拡大}
\begin{defn} (代数閉体・代数閉包)\\
\ben
\item $K$が代数閉体(algebraic closed field)であるとは, 任意の定数でない$f(x) \in K[x]$が$K$に根をもつことをいう.
\item 体$K$を含む代数閉体$L$を$K$の代数閉包(algebraic closure)という.
\een
\end{defn}

\begin{rem}
任意の代数閉体$K$は無限集合である.標数$0$なら$\mb{Q}\subset K$だから明らか.標数$p$かつ$|K| = n$だと仮定すると,$n$より大きい$p$で割れない整数$N$を取って $K$係数の$N$次多項式$x^{N} - 1$は分離多項式であり, $K$が代数閉体であるという仮定から1次式に分解されるが, $|K|=n$なのでそのとき同じ根が二つ以上現れてしまい矛盾である.よって$K$は無限集合である.
\end{rem}

\begin{theo}
任意の体$K$は代数閉包を同型を除いてただひとつ持つ.したがって$K$の代数閉包を$\ovl{K}$と書く.
\end{theo}

\begin{egprob}
$L/K$を代数拡大, $\alpha_1,\dots, \alpha_n\in L$を$K$上代数的な元とするとき, $K[\alpha_1,\dots, \alpha_n] = K(\alpha_1,\dots, \alpha_n)$であることを示せ.
\end{egprob}
\begin{egans}
まず両辺の集合の意味を復習しておこう. 左辺は, $\set{\alpha_i}$で生成される有限生成$K$代数であり, 多項式の代入した元の全体になる. 右辺は, 有理関数(で分母が0にならないもの)を代入した元で生成される体である. \par
左辺は実は体である. [証明] $\alpha_i$が$K$上代数的という仮定から, $K[\alpha_1,\dots, \alpha_n]\supset K$が整拡大であって, $K$が体であることから環論的に従う. [終] よって, どのような$f(\alpha_1,\dots, \alpha_n)$の逆元も存在し, $K(\alpha_1,\dots, \alpha_n)\subset K[\alpha_1,\dots, \alpha_n]$となる. 逆の包含は自明.\qed
\end{egans}

\begin{prac}
体拡大$L/K$について, 任意の$\alpha \in L$に対して$K[\alpha] = K(\alpha)$であることを示せ.
\end{prac}


\subsection{分離拡大}
\begin{prop} (\tf{分離拡大まとめ})
\ben
\item \tf{有限次分離拡大は単拡大である. }
\een
\end{prop}
\prf
(1): $K$の2元生成分離拡大体$L=K(\alpha_1,\alpha_2)$が$K(\alpha)$と書けるような$\alpha=\alpha_1+c\alpha_2$を発見できれば帰納的に示せる. まず, $K$が有限体であった場合は乗法群が生成元であることから明らかなので, $K$は無限体であるとする. $[L:K] = n$とする. 分離性より$|\mr{Hom}_{K}(L,\ovl{K})| = n$である. $K(\alpha) = L$であることを示すには, $[K(\alpha):K] = |\mr{Hom}_{K}(K(\alpha), \ovl{K})| = n$であることを示せばよいから, $\alpha$が満たすべき条件というのは次である:
\lefba{
任意の相異なる$\phi_i,\phi_j\in \mr{Hom}_{K}(L,\ovl{K})$に対して$\phi_i(\alpha) \neq \phi_j(\alpha)$である.
}
このような$c,d$が発見できればよい. そこで, 次の多項式を考える.
\[F(x) = \prod_{i\neq j}(\phi_i(\alpha_1) + x\phi_i(\alpha_2) - \phi_j(\alpha_1) - x\phi_j(\alpha_2)  )\in L[x]\]
さて, $L=K(\alpha_1,\alpha_2)$なので$\phi\in \mr{Hom}_{K}(L,\ovl{K})$は$\alpha_1,\alpha_2$の写し方で決定されていることに注意すると, 上の$F$の各因子は零多項式ではない. したがって$F(x)$は0多項式ではない. したがって, \tf{$K$は無限体なので}ある$x=c$が存在して$F(c) \neq 0$である. ゆえに, このように$c$を取れば$\phi_i(\alpha)\neq \phi_j(\alpha)$が言えるので$L=K(\alpha)$が成り立つ.\qed


\subsection{$\spadesuit$純非分離拡大}

\subsection{正規拡大}
\begin{prop} (\tf{正規拡大まとめ})
\ben
\item $L/K$が正規拡大のとき, 任意の$\sigma\in\mr{Hom}_{K}(L,\ovl{K})$に対し$\sigma(L)\subset L$である. よって$\mr{Hom}_{K}(L,\ovl{K}) = \mr{Hom}_{K}(L,L)$.
\item $L=K(\alpha_1,\dots, \alpha_n)$と書けるとき, これが正規拡大なら$\alpha_i$の$K$共役元はまた$L$に含まれる. 逆にすべての$\alpha_i$の$K$共役元を$L$が持つならば, $L/K$は正規拡大である.
\een
\end{prop}

\begin{tips}
\ben
\item 2次の体拡大は\tf{標数2の場合を除き}正規拡大である.これは群論からすれば「指数2の部分群は正規部分群である」ということに対応する.
\een
\end{tips}

\begin{egprob}
$L/M,M/K$が正規拡大のとき$L/K$もまた正規拡大か?
\end{egprob}
\begin{egans}
No. $L=\mb{Q}(\sqrt[4]{2})$, $M=\mb{Q}(\sqrt{2})$, $K=\mb{Q}$とせよ.
\end{egans}

\begin{egprob} (2020代数学II, 11/10練習問題)
次の拡大体は正規拡大か? ただし, $n\geq 3$とし, $\zeta = \zeta_n$は1の原始$n$乗根とする.
\ben
\item $\mb{Q}(\sqrt{2},\sqrt{3},\sqrt{5})/\mb{Q}$
\item $\mb{Q}(\zeta)/\mb{Q}$
\item $\mb{Q}(\sqrt[n]{2})/\mb{Q}$
\item $\mb{Q}(\sqrt[n]{2},\zeta)/\mb{Q}$
\item $\mb{Q}(\sqrt[n]{2},\zeta)/\mb{Q}(\zeta)$
\item $p$を素数, $\mb{F}_{p}(X)$を1変数有理関数体とし, $\mb{F}_{p}(X)/\mb{F}_{p}(X^p)$.
\een
\end{egprob}
\begin{egans}
(1):○ (2):○ (3):$\zeta$の共役元は$\mb{R}$にないものもあるので× (4):○ (5):正規拡大から下の体を挙げても正規なので○. (6): $X$の$\mb{F}_{p}(X^p)$上の多項式$T^p - X^p$の零点であるが, これを分解すると$(T-X)^p$なので, 共役元は$T=X$のみであると分かる. したがって正規拡大.
\end{egans}



\subsection{Galois拡大}
Galois拡大とは, 分離拡大かつ正規拡大であるような体拡大である. 分離拡大$L/K$に対し, $L$の$K$上の\tf{Galois閉包}$\witi{L}$とは, $L$の$K$上の共役元をすべて添加した体のことをいう. $L$の$K$代数としての自己同形全体の群$\mr{Aut}_{K}^{\mr{alg}}L$を一般に$\gal{(L/K)}$と表し, $L/K$のガロア群と呼ぶ. \par
\tbox{
\begin{prop} (\tf{ガロア群まとめ}) $L/K$を有限次ガロア拡大とする.
\ben
\item $\gal{(L/K)}$の位数は$[L:K]$である.
\item $M$が中間体のとき, $L/M$もガロア拡大である.
\item $L=K(\alpha)$と書けるとする. $\sigma\in \gal{(L/K)}$は$\sigma(\alpha)$で決定され, $\sigma(\alpha)$は$\alpha$の$K$共役元に移る. したがって, $\alpha$が$d$次の$K$係数多項式の根になることが分かっている場合, $\sigma(\alpha)$としてあり得るのは$d$通りであり, $[L:K]\leq d$が従う.
\een
\end{prop}
}{title=ガロア群まとめ}
\prf
(1): 分離性より$|\mr{Hom}_{K}(L,\ovl{K})| = [L:K]$. 正規性より$\mr{Hom}_{K}(L,\ovl{K}) = \mr{Hom}(L,L) = \mr{Aut}_{K}L = \gal{(L/K)}$だから.
\qed

\begin{rem} ここで, 一般の$K$準同型に関する簡単ですが重要な命題を述べます. $S$が$K$代数$A$の生成系, すなわち$A=K[S]$であるとします. $K$代数の準同型$\phi:A\to B$は, $\phi(s)$ ($s\in S$)の値によって一意に決まります. とくに, $A=B=L$で$L$が単拡大$L=K(\alpha)$の場合, $\phi(\alpha)$の値を見れば$\phi$は一意に決まります. この事実は, ガロア群の位数を評価するときなどに用いられます.
\end{rem}

$L/M$はガロア拡大になるのに対し, 次には注意しよう.
\tbox{
\begin{egprob}
$L/K$をガロア拡大とし, $M$を中間体とする. $M/K$はガロア拡大であるか?
\end{egprob}
}{}
\begin{egans}
一般にならない. たとえば$L=\mb{Q}(\omega,\sqrt[3]{2})$, $M=\mb{Q}(\sqrt[3]{2})$とすると$M/\mb{Q}$は正規拡大ではない.
\end{egans}



\begin{egprob}
次の体$L$の$\mb{Q}$上のガロア閉包$F$を求め, $\gal{(F/\mb{Q})}$を決定せよ: $L=\mb{Q}(\sqrt[3]{3})$/
\end{egprob}
\begin{egans}
$L=\mb{Q}(\sqrt[3]{3})$のとき, $\mr{Irr}(\sqrt[3]{3},\mb{Q};x) = x^3 - 3$であり, 共役元は$\sqrt[3]{3}\omega^{k}$ ($k=0,1,2$)である. よって, $F = \mb{Q}(\sqrt[3]{3}, \sqrt[3]{3}\omega, \sqrt[3]{3}\omega^2) = \mb{Q}(\sqrt[3]{3},\omega)$. これの拡大次数は6の倍数であること, および6以下であることが容易に分かるので拡大次数は6である. よって$|\gal{(L/K)}| = 6$.  $\sigma\in \gal{(L/K)}$は集合$\set{\sqrt[3]{3}, \sqrt[3]{3}\omega, \sqrt[3]{3}\omega^2}$の置換を引き起こし, 逆に, この集合が$L$を生成するのでその置換の情報から$\sigma$は一意に決定される. 置換は6通りあり, ガロア群の位数も6通りなので, すべての置換を実現できる. すなわち, 次のような$\sigma, \tau\in \gal{(L/K)}$が存在する:
\[\sigma(\sqrt[3]{3}) = \sqrt[3]{3}\omega, \sigma(\sqrt[3]{3}\omega) =\sqrt[3]{3}\]
\[\tau(\sqrt[3]{3}\omega^k) = \sqrt[3]{3}\omega^{k+1}\]
これを生成系$\set{\sqrt[3]{3},\omega}$の行き先から表現する式に書き換えると
\[\sigma(\sqrt[3]{3}) = \sqrt[3]{3}\omega,\quad \sigma(\omega) = \omega^2\]
\[ \tau(\sqrt[3]{3}) = \sqrt[3]{3}\omega , \tau(\omega) = \tau(\omega) = \omega \]
である. $\sigma$は互換に対応しており, $\tau$は3-cycleに対応している. $\gal{(L/K)}$はこの$\tau,\sigma$で生成される群である.
\end{egans}





\newpage
\section{Galois理論の基礎}
\subsection{ガロア拡大とArtinの定理}

\tbox{
\begin{theo}
$G$を有限群, $L$を体, $\rho:G\to \mr{Aut}L$を単射準同型とし,
\[K := L^{G} := \{ \alpha \in L \mid \forall g\in G,  \rho(g) (\alpha) = \alpha \}\]
とおく.このとき, $L/K$は有限次ガロア拡大であり,
\[\bolm{\gal(L/L^{G}) \cong G}\]
である.\
\end{theo}
}{title=Artinの定理}
\begin{point}
この定理の証明のめんどくさい部分は「有限次ガロア拡大であることを示す部分」であり, そこを認めれば実は同型は位数比較で示せる.だから, 特定の状況下で有限次ガロア拡大であることが分かっていたら, かなり使い勝手がよいものである.(It先生曰く一番重要な定理である)
\end{point}
\prf \tf{有限次ガロア拡大であることを認めて最後の同型を示す.} 有限次ガロア拡大であることは\cite{aoyukie}などを見よ. \par
次の不等式に注意しよう.
\[|G| \leq |\mr{Aut}_{L^G}L| = |\mr{Gal}(L/L^G)| \leq |\mr{Hom}_{L^G}(L,\ovl{K})| = [L:L^G]  \]
最初の不等号は$\rho$の単射性による. よって$[L:L^G] \leq |G|$を示せばよくなる. $L/L^G$を有限次ガロア拡大と認めたので, とくに単拡大である. $L=L^G(\alpha)$と書ける. \ao{$[L:L^{G}]\leq |G|$を示すには, $\alpha$が$L^G$係数の$|G|$次以下の多項式の根になればよい. } $\set{\rho(g)(\alpha) \mid g\in G } = \set{\alpha_1,\dots, \alpha_n}$とおき, $i\neq j$なら$\alpha_i\neq \alpha_j$であるようにする(よって\bolm{n\leq |G|}である). このとき$\alpha_1,\dots,\alpha_m$の第$i$次対称式$s_i$は$\rho:G\to \mr{Aut}L$の作用で不変であり,
\[f(t) = \prod_{i=1}^{m} (t-\alpha_i) = t^n - s_1t^{n-1} + \dots + (-1)^{m} s_m\]
は$L^G[t]$の元で, $y$を根に持つ$L^G$係数の$n$次多項式が得られたから$[L^G(\alpha):L^G] \leq n\leq |G|$である.  $\rho$の像は$\mr{Aut}_{L^G}L$に含まれるから, $\rho:G\to \mr{Aut}_{L^G}L$が全単射となり同型が従う.

\begin{comment}
$n=[L:L^G]$とする.$L^{G}$の定義から, $\rho(G) \subset \mr{Aut}_{L^G}L = \gal(L/L^G)$であり, $\rho$は単射なので$|G|\leq n$である.\par
逆の不等号を示す.$|G| < n$であるとする. $t> |G|$で, $x_1,\cdots, x_t\in L$を$L^G$上一次独立とするなら, $[L^G(x_1,\cdots, x_t):L^G] > |G|$だが, $L^G(x_1,\cdots, x_t)$は単拡大$L^G(y)$と書け, $y$の最小多項式は$|G|$次より大であるということになる.\par
$y$の$G$の作用による軌道の元を$\{ y=y_1,\cdots, y_r\}$とするとき,
\ben
\item $G$軌道が$r$個の元しか作らないのだから当然$r\leq |G| $である.
\item $s_i$を$y_j$の基本対称式とすれば
\[F(T) = \sum_{i=0}^{r} (-1)^i s_i T \in L^G[T]\]
は$\{ y=y_1,\cdots, y_r \}$を根に持つので, \aka{最小性から$r > |G|$である.}
\een
これは矛盾である.よって$|G| \geq n$が示され, $|G| = n$である.よって, $\rho:G\to \gal{(L/L^G)}$は群同型である.   \qed
\end{comment}



\begin{eg}
$G=\maf{S}_{n}$は有理関数体$L=k(x_1,\cdots, x_n)$に置換として忠実に作用する。\textcolor{blue}{よって命題より
\[\mr{Gal}(L/L^{G}) \cong \maf{S}_{n}\]
だから$[L:L^{G}] = n!$である。}

$x_1,\cdots, x_n$の基本対称式を$s_1,\cdots, s_n$とする。不変体$L^{G}$が何であるかを調べよう。$k(s_1,\cdots, s_n)$はいかにも不変体らしく見える。実際そうである。$F=k(s_1,\cdots, s_n)$としよう。$F\subset L^{G}$は明らかである。よって$[L:F] \geq n!$である。逆の包含関係を示すためには, 拡大次数が$\leq n!$であると言えればよい。(このような使い方が典型的)

まず$L/F$はガロア拡大である。なぜなら, $L= F(x_1,\cdots, x_n)$であり, $T^{n} - s_1T + \cdots + (-1)^{n}s_n$は$x_i$をすべて根にもつから正規であり, この多項式$\in F[T]$は重根を持たないから分離的である。 係数は$G$で不変であるから, $\sigma\in \mr{Gal}(L/F)$は$f(t)$の根の置換を起こす。\ao{よって準同型$G\to \maf{S}_n$が得られる。}逆にこの置換を一つ定めることで$\sigma$は完全に決定されるので準同型は単射。よって$|\mr{Gal}(L/F)| = [L:F]\leq n!$である. 以上より$F=L^{G}$, すなわち次が言えたことになる.

\textcolor{blue}{有理関数のうち, 置換を施して変わらないものはいずれも基本対称式の有理関数として表すことができる。}
\end{eg}




\begin{egprob} (東工大R2午後2)
$k$を標数$p>0$の体とし, $n\in\mb{N}$とする. $n$変数有理関数体$L=k(t_1,\dots, t_n)$を考える. $V$により, $t_1,\dots, t_n$の生成する$L$の$n$次元部分$\mb{F}_{p}$ベクトル空間を表す. $\mr{GL}_n(\mb{F}_{p})$を$L$に次のように作用させる: $g=(g_{ij})\in \mr{GL}_{n}(\mb{F}_p)$と$f=f(t_1,\dots, t_n)\in L$に対し,
\[(gf)(t_1,\dots, t_n) = f(g(t_1), \dots g(t_n)),\quad g(t_j) = \sum_{i=1}^{n} g_{ij}t_i.\]
さらに$F(X) = \prod_{v\in V} (X-v)$とおく.
\ben
\item $F(X) = X^{p^n} + s_1 X^{p^{n-1}} + \dots + s_{n-1}X^p + s_n X$, $s_i \in L^{\mr{GL}_{n}(\mb{F}_p)}$と表せることを証明せよ. ここで$L^{\mr{GL}_{n}(\mb{F}_{p})}$は$L$の$\mr{GL}_{n}(\mb{F}_p)$による固定部分体を表す.
\item $L$の部分体$K=k(s_1,\dots, s_n)$を考える. $L/K$は有限次ガロア拡大であり, そのガロア群は$\mr{GL}_{n}(\mb{F}_{p})$と同型であることを証明せよ.
\een
\end{egprob}

\begin{egans}
(1): 帰納的に考える. $n=1$のときは$F(X) = X^{p} - Xt_1^{p-1}$なのでよい.  \par
$F(X)$の係数が$L^{\mr{GL}_{n}(\mb{F}_{p})}$に属すことは, $v\in V$たちの基本対称式で表されていることから明らかである.
$W=\mb{F}_{p}t_1 \oplus \dots \oplus \mb{F}_{p}t_{n-1}$とするとき, \[F(X) = \prod_{w\in W}\prod_{m=0}^{p-1} (X-w-mt_n)\]
である. $X-mt_n$を$X_m$と書いて一文字と見ると,
\[\prod_{w\in W} (X_m - w) = X_m^{p^{n-1}} + s_1X_m^{p^{n-2}} + \dots + s_{n-1} \]
と表せる. これを$G(X_m)$と書くと, 指数が$p$の冪の部分しかないので, $G(X_m) = G(X) - mG(t_n)$が分かる. よって,
\[F(X) = \prod_{m=0}^{p-1} (G(X) - mG(t_n)) = G(X^{p}) - G(X) G(t_n)^{p-1} \]
となり, 明らかに$p$冪次以外の部分は消失している. \\
(2): (本問はArtinの定理を用いるが, 答案に書くための流れでいきたい)\\
まず$t_i$は$F(X)$の根であり, $t_i$の共役元は$F(X)$の根の中にあることは分かるので, $L/L^{\mr{GL}_{n}(\mb{F}_p)}$が正規拡大であることは分かる. そして, $F'(X) = s_n \neq 0$なので分離拡大である(この拡大では下の体が完全体ではないのでこの確認は必須). よって有限次ガロア拡大である. \\
\fbox{$\gal{(L/L^{\mr{GL}_{n}(\mb{F}_p)})} \cong \mr{GL}_{n}(\mb{F}_p)$} $G=\mr{GL}_{n}(\mb{F}_p)$とおく. 問題文によって忠実作用$\rho:G\to \mr{Aut}_{L^{G}}L$が与えられる. 次に注意する.
\[|G| \leq |\mr{Aut}_{L^G}L| \leq |\mr{Hom}_{L^G}(L,\ovl{L^G})| = [L:L^G]\]
$L/L^G$は有限次分離拡大だから$L=L^G(\alpha)$と表せる. $\rho(g)(\alpha)$という形の元のうち異なるものを$\alpha=\alpha_1,\dots, \alpha_m$とすると$m\leq |G|$である. そして, $\alpha_1,\dots, \alpha_m$の基本対称式$l_1,\dots, l_m\in L^G$を用いて
\[T^m - l_1T^{m-1} + \dots + (-1)^{m}l_m\]
は$\alpha_1,\dots, \alpha_m$を根に持つので$[L^G(\alpha):L^G]\leq m \leq |G|$である. よって$\rho$は全単射だから同型. ゆえに$\gal{(L/L^G)}\cong G$. \\
\fbox{$k(s_1,\dots, s_n) = L^{G}$} $\subset$は明らかである. よって $L/k(s_1,\dots,s_n)$もガロア拡大であり, $[L:k(s_1,\dots, s_n)] \leq |G|$が示せれば拡大次数を比較して$L^G=k(s_1,\dots, s_n)$が従う.
$\sigma\in \gal{(L/k(s_1,\dots, s_n))}$は$\sigma(t_i)$の移り方で決まる. $\sigma(t_i)$の写し先としてあり得るのは$V$の元である. しかし, $\sigma$は同型写像で$V$の元は$t_1,\dots, t_n$で生成されているので, $\sigma(t_1),\dots, \sigma(t_n)$の移り方は, $\sigma$を$V$に制限して置換になるようにするべきである. このため, $\sigma(t_1),\dots, \sigma(t_n)$は$V$の一次独立系でなければならない. よって$\sigma$にただひとつの$G$の元が(表現行列として)対応するべきであるから, 単射準同型$\gal{(L/k(s_1,\dots, s_n))}\hookrightarrow G$がある.よって$|\gal{(L/k(s_1,\dots, s_n))}|\leq |G| = [L:L^G]$であるから, これと$k(s_1,\dots, s_n)\subset L^G$より$L^G = k(s_1,\dots, s_n)$である.\qed
\end{egans}




\subsection{Galois理論の基本定理}
Galois群の部分群と拡大体の中間体が1:1に対応するということの正確な主張は以下の通りである.
\tbox{
\begin{theo} $L/K$を有限次ガロア拡大とする. $\mb{M}$を$L/K$の中間体の集合, $\mb{H}$を$\gal{(L/K)}$の部分群の集合とするとき, 次の(1)-(3)が成り立つ.
\ben
\item 次の二つの写像は互いの逆写像である:
\bealn{
\mb{M}\ni M  &&\mapsto && H(M) \in \mb{H}\\
\mb{M} \ni M_{H} && \mapsfrom && H\in \mb{H}
}
ただし, $H(M)$は$M$を固定するガロア群の元全体$\set{g\in \gal{(L/K)} \mid \forall x\in M, gx = x}$であり, $M_H$は$H$により固定される中間体$\set{x\in L \mid \forall g\in H \mid gx = x}$である(不変体という).
\item $M_1,M_2\in\mb{M}$が$H_1,H_2\in\mb{H}$に対応するとき, 『$\cap$とSpanで双対的かつ順序を反対にする』対応になる.
\bealn{
M_1\subset M_2 &&\iff && H_1\supset H_2&\quad (\tf{反変性}) \\
M_1\cdot M_2 && \leftrightarrow && H_1\cap H_2
 & \\
 M_1 \cap M_2 && \leftrightarrow && \lan H_1,H_2 \ran
 }
 \item $M\in \mb{M}$が$H\in \mb{H}$と対応し, $\sigma \in \gal{(L/K)}$ならば,
 \[\sigma(M) \leftrightarrow \sigma H \sigma^{-1}\]
 と対応する. さらに,
 \[\bolm{ M/K がガロア拡大 \iff H が \gal{(L/K)}で正規}\]
 である. また, $H$が正規部分群であるとき, $\gal{(L/K)}$の元を$M$に制限することにより
 \[\bolm{\gal{(M/K)} \cong \gal{(L/K)}/H}\]
 である.
\een

\end{theo}
}{title=ガロアの基本定理}
\prf
\step{1.1}{$M\mapsto H(M) \mapsto M_{H(M)}$は恒等写像}
定義から$M_{H(M)}$は$H(M) = \gal{(L/M)}$の全ての元で固定される体であるから$M_{H(M)} \supset M$が従う. $[L:M] = n$としよう. $|H(M)| = n$である. 包含射$\rho:\gal{(L/M)}\to \mr{Aut}L$の作用に関する$L$の固定体が$L^{\gal{(L/M)}} = M_{H(M)}$であるから, アルティンの定理により$\gal{(L/M_{H(M)})} \cong \gal{(L/M)}$であり(\tf{ここがポイント}), とくに位数が等しいので$[L:M_{H(M)}] = n$である. $M_{H(M)}\supset M$なので, 等号が成り立つ. \\
\step{1.2}{$H\mapsto M_H\mapsto H(M_H)$は恒等写像}
定義により$H(M_H) = \gal{(L/M_H)}$であり, $H$の元は$M_H$を固定するので$H \subset H(M_H)$である. $\rho:H\to \mr{Aut}L$により$L$に忠実に作用し, $L^H := M_H$なのでアルティンの定理から$\gal{(L/M_H)} = H(M_H) \cong H$(ここがポイント). とくに$H$と$H(M_H)$の位数は等しいので$H=H(M_H)$.
(2): \\
(3):
\qed



\begin{egprob}
$K/\mb{Q}$を拡大で, ある$n$により$K\subset \mb{Q}(\zeta_n)$であるとするとき, $K/\mb{Q}$はガロア拡大であることを示せ. また, これを用いて $K=\mb{Q}(\tan{\dfrac{2\pi}{n}})$が$\mb{Q}$上ガロア拡大であることを示せ.
\end{egprob}
\begin{egans}
$\mb{Q}(\zeta_n)/\mb{Q}$はアーベル拡大であるから, 中間体に対応する部分群は必ず正規部分群であり, そのため$K/\mb{Q}$がガロア拡大になる. $K=\mb{Q}(\tan{\dfrac{2\pi}{n}})$の場合, $K\subset \mb{Q}(\zeta_{4n})$なのでよい($\zeta_{4n}^{n} = i$に注意).
\end{egans}


次は過去問ですが, ここで用いる議論はほかの問題にも応用できます.

\begin{egprob} (H28専門科目1) \label{prob:H28senmon-1}
$\mb{Q}$係数既約多項式$f(x) = x^3 + ax+b$を考え, $K\subset \mb{C}$を$f(x)$の最小分解体とする. $a>0$のとき, $K/\mb{Q}$のガロア群が3次対称群と同型であることを示せ.
\end{egprob}
\begin{egans}
$f'(x) = 3x^2+a>0$なので中間値の定理より実根$r$がただひとつ存在する. 虚根を$\alpha$とすると$\ovl{\alpha}, \alpha , r$が根のすべてになる. $r+\alpha + \ovl{\alpha} = 0$なので$\ovl{\alpha} \in \mb{Q}(\alpha,r)$. よって$K=\mb{Q}(\alpha,r)$である. これが$\mb{Q}$上の6次ガロア拡大であることは容易に分かる. ガロア群を$G$とする. $\sigma\in G$は$\sigma(r)$と$\sigma(\alpha)$によって決まる. これは$X=\set{r,\alpha,\ovl{\alpha}}$の置換によって決まるということと同じである. よって$G$は$X$に作用し, $G\to \mr{Aut}(X)\cong \maf{S}_{3}$という群準同型を得るが, これは位数が同じで単射なので全単射でもある. よって同型である. \qed
\end{egans}

\begin{prac}
$a=0$の場合にも同様の議論が成り立つか確かめよ.
\end{prac}

\begin{egprob} (仮想面接問題)
先の例題で$a<0$の場合に$S_3$と同型にならない例をあげてください. 難しければ$f(x)=X^3+aX^2+bX+c$の最小分解体からでもよいです.
\end{egprob}

\begin{egans}
実根が3つあるということが起こりうるので, それを狙う. $t_k=\zeta_7^k + \zeta_7^{-k}$とし, $t_1$の最小多項式を求める. $2\in (\cyc{7})^{\ast}$が位数3であることを利用して,
\[(X-t_1)(X-t_2)(X-t_4)\]
とする.
\bealn{
t_1+t_2+t_4 &= -1\\
t_1t_2 + t_2t_4 + t_4t_1 = -2 \\
t_1t_2t_4 &= 1
}
と計算でき, $X^3 + X^2 - 2X - 1$を得る. これの最小分解体は$\mb{Q}(t_1)$である(普通に計算するか, チェビシェフ多項式などを考えよ). よってこの拡大は3次であり, ガロア群は$\cyc{3}$である.
\end{egans}
\begin{egans} 案外, $S_3$にならない例というのは限られている.  ここでの例は「三次多項式のガロア群」の節を参照すると具合がわかる. 実は$f(x) = x^3-3x-1$の最小分解体が$\cyc{3}$である. 解答は\ref{egprob:gal=z/3z_example}にて.
\end{egans}










\subsection{Galois拡大の推進定理}
合成体のGalois群は推進定理を考えると良いことがある.
\tbox{
\begin{theo}
$L = M_1\cdot M_2$, $M_1\cap M_2 = K$とする.
\ben
\item $M/K$が有限次ガロア拡大ならば, $L/N$も有限次ガロア拡大で,
\[\gal{(L/N)}\cong \gal{(M/K)}\]
\item $M/K, N/K$が有限次ガロア拡大ならば$L/K$も有限次ガロア拡大で,
\[\gal{(L/K)}\cong \gal{(M/K)}\times \gal{(N/K)}\]
\een
\end{theo}
}{title=Galois拡大の推進定理}
\begin{point}
\ben
\item 「平行四辺形」の向かい合う辺のうち下にあるほうがFin.Galならば, もう片方も同じガロア群を持つFin.Galである.
\item 下の2辺がFin.Galならば, 下から上までのガロア群は直積.
\een
\end{point}
\begin{cor} (簡易版1)
上の設定で, $L/K$, $M/K$が有限次ガロア拡大ならば, $L/N$も有限次ガロア拡大で
\[\gal{(L/N)}\cong \gal{(M/K)}\]
であり, $N/K$も有限次ガロア拡大ならば
\[\gal{(L/K)}\cong \gal{(M/K)}\times \gal{(N/K)}\]
\end{cor}

\begin{eg}
有限次ガロア拡大でない拡大がある場合の使用例を一つ考える. $\mb{Q}\subset \mb{R}\subset \mb{C}$という体拡大を考え, $K/\mb{Q}$を有限次ガロア拡大で, $K\neq K\cap \mb{R}$であるとする. このとき$K\cdot \mb{R} = \mb{C}$に注意する. もし$K\cap \mb{R} = \mb{Q}$であれば, 定理の(1)によって
\[\gal{(K/\mb{Q})}\cong \gal{(\mb{C}/\mb{R})} \cong \mb{Z}/2\mb{Z}\]
である. よって$K$は虚2次体$\mb{Q}(\sqrt{-d})$である. $K\cap \mb{R}$が分からなくても, $K/K\cap \mb{R}$に上を適用すれば\tf{これが2次拡大であると分かる.}
\end{eg}
\begin{egprob} $x^3+2$の$\mb{Q}$上の最小分解体$K$に対し, $K\cap \mb{R}$を求めよ.
\end{egprob}
\begin{egans}
$K=\mb{Q}(\sqrt[3]{2}, \zeta_6)$が用意に分かる. $K$は$K\cap \mb{R}$上2次なので$[K\cap \mb{R}:\mb{Q}] = 3$である. 明らかに$K\cap \mb{R}\supset \mb{Q}(\sqrt[3]{2})$なので$K\cap \mb{R} = \mb{Q}(\sqrt[3]{2})$.
\end{egans}
\begin{prac}
$K/\mb{Q}$が奇数次ガロア拡大のとき$K\subset \mb{R}$を推進定理で示せ\footnote{推進定理なしなら, 複素共役という自己同形をがあるのを考えれば簡単. }.
\end{prac}







\subsection{Kummer理論}
\begin{assume} \label{assume:kummer}
$n$を固定した2以上の正の整数とする. 体$K$は, 標数$p$が$p=0$あるいは$n$を割り切らない素数$p$であるとして,  1の原始$n$乗根$\zeta_n$を含むものであるとする.
\end{assume}
\begin{eg}
\ben
\item $n=2$なら, $\mb{Q}$や$\mb{F}_{p^{k}}$ ($p$は奇素数) は条件を満たす.$a_i\in \mb{Z}$として, \textcolor{blue}{$\mb{Q}(\sqrt{a_1}, \sqrt{a_2}, \cdots, \sqrt{a_k})$のような拡大体も条件を満たす.}
\item $n=3$なら, $\mb{Q}(\omega) = \mb{Q}(\sqrt{-3})$や$\mb{F}_{p}$ ($p=6k+1$型) は条件を満たす.
\een
\end{eg}

\tbox{
\begin{theo} $K$はある$n\geq 2$に対して仮定\ref{assume:kummer}を満たす体とする.$R$は部分群$K^{\times} \supset R \supset (K^{\times})^{n}$で, $R/(K^{\times})^{n}$が有限であるとする.このとき$K(\sqrt[n]{R})/K$はGalois拡大であって, そのGalois群に対して
\[\gal{(K(\sqrt[n]{R})/K)} \cong \hom{{R/(K^{\times})^{n}}, \mb{C}^{1}}\]
が成り立つ.ただし, 右辺は$R/(K^{\times})^{n}$から$\mb{C}^{1}= \{ z\in \mb{C} \mid |z| = 1 \}$ への群準同型全体の集合に\tf{通常の積の演算}を入れたものである\footnote{有限Abel群$G$に対して$\hom{G,\mb{C}^{1}}$を$G^{\ast}$とも書く本もある(\cite{aoyukie}).また, $G^{\ast} = \hom{G,\mb{C}^{\times}}$である.}.
\end{theo}
}{title=\tf{Kummer拡大のGalois群}}
すなわち, $K$に有限的な条件を守りながら冪根の集合$R$を添加した体$K(\sqrt[n]{R})$の$K$上のガロア群は, 群$R/(K^{\times})^n$の\tf{双対群}に同型である. \\
\prf
方針を述べる. $\mu_n$を$\mb{C}_1$の$1$の$n$乗根の集合とする.
\ben
\item $\Phi:\gal{(L/K)}\times R\to \mu_n$を
\[\bolm{\Phi(\sigma, a) = \dfrac{\sigma(\sqrt[n]{a})}{\sqrt[n]{a}}}\]
と定めると, これは$\sqrt[n]{a}$の選択に依存しない.
\item $\Phi$はPerfect Pairing $\gal{(L/K)}\times R/(K^{\times})^n \to \mu_n$を引き起こす. すなわち, この写像は非退化$\mb{Z}$双線型写像である.
\item 一般にPerfect Pairing$\Phi:G\times H \to \mb{C}_1$があｒとき, $H\cong G^{\ast}, G\cong H^{\ast}$である,
\een
\qed


\begin{eg}
素数を小さい順に並べた数列を$p_1,p_2,\cdots$とする.$L = \mb{Q}(\sqrt{p_1},\sqrt{p_2},\cdots, \sqrt{p_k})$の拡大次数は$2^{k}$であると予想できる.実際正しい.初等整数論的に解けるかもしれない.だがこれはよく見るとKummer拡大であることが分かるので,  Kummer理論を使って拡大次数を求めることができるのである.そのためには$L$が$\mb{Q}(\sqrt[2]{R})$と思えるように部分群$\mb{Q}^{\times} \supset R \supset (\mb{Q}^{\times})^{2}$を取らなければならない. ここれは, $R$を$p_1,\dots, p_k$と$(\mb{Q}^{\times})^2$で生成される$\mb{Q}^{\times}$の部分群とする. $R/(\mb{Q}^{\times})^2 \cong (\mb{Z}/2\mb{Z})^k$が示せる. これの双対群もまた$(\mb{Z}/2\mb{Z})^k$である.
\end{eg}

\begin{eg}
$l$を素数, $\mu\in \mb{Z}$を整数の$l$乗ではないとする.$K=\mb{Q}(\sqrt[l]{\mu})$とするとき, $[K:\mb{Q}] = l$であることを示そう.
\end{eg}



\begin{egprob} (\cite{momoyukie}演習2.2.2)\\
$K=\mb{Q}(\sqrt[5]{2},\sqrt[5]{3})$, $\zeta = \exp{(2\pi \sqrt{-1}/5)}$, $L=K(\zeta)$とする.
\ben
\item $L/\mb{Q}(\zeta)$はガロア拡大で$\gal{L/\mb{Q}(\zeta)} \cong (\mb{Z}/5\mb{Z})^2$であることを証明せよ.
\item $[K:\mb{Q}] = 25$であることを証明せよ.
\een
\end{egprob}
\begin{egans}
(1): $K_0 = \mb{Q}(\zeta)$とする. $R$を$(K_0^{\times})^5$と$2,3$で生成される群とする. $K_0$を1の5乗根を含む体である. $K=K_0(\sqrt[5]{R})$と表せるので, Kummer理論により
\[\gal{(K_0(\sqrt[5]{R})/K_0)} \cong (R/(K_0^{\times})^5)^{\ast}\]
である. $f:\mb{Z}^2 \to K_0^{\times}/(K_0^{\times})^5$を$f(a,b) = [2^a3^b]$とする. $\ker{f} = (5\mb{Z})^2$であることを示そう. もし$2^a3^b\in (K_0^{\times})^5$に属すなら, $2^a3^b = \alpha^5$なる$\alpha\in K_0$が存在する. これにより
\[N_{K_0/\mb{Q}}(2^a3^b)= 2^{4a}3^{4b} = N_{K_0/\mb{Q}}(\alpha^5) = (N_{K_0/\mb{Q}}(\alpha))^5 \]
であり, $N_{K_0/\mb{Q}}(\alpha) \in \mb{Q}$なので$a,b\in 5\mb{Z}$. 逆にこのときに$\ker{f}$に属すのは明らか. よって$\ker{f} - (5\mb{Z})^2$なので, 準同型定理より$\im{f} \cong (\mb{Z}/5\mb{Z})^{2}$である. \qed \\
(2): $K\cap K_0 \subset \mb{R} \cap \mb{Q}(\zeta) = \mb{Q}$なので$K\cap K_0 = \mb{Q}$. よって推新定理より$K_0/\mb{Q}$と同じく, $L/K$は4次拡大. よって(1)$[L:\mb{Q}] = 4[K:\mb{Q}]$である. 一方で(1)から$[L:\mb{Q}] = [L:K_0][K_0:\mb{Q}] = 25\times 4  =100$だから, $[K:\mb{Q}] = 25$が従う.\qed
\end{egans}





\newpage


\subsubsection{演習問題}
\subsubsection{Level A}
\begin{prob} (20代数学IIレポート)\\
$f(T)\in \mb{Q}[T]$を既約な3次多項式で, 3つの実根をもつものとする.$\alpha\in \ovl{Q}$を$f(T)$の根とする.このとき, 次の条件をみたす$\beta\in /mb{Q}(\alpha)$は存在しないことを示せ.
\[\beta^{3}\in \mb{Q}, \mb{Q}(\alpha) = \mb{Q}(\beta)\]
\end{prob}

\subsubsection{Hint・Answer}
{\small
\begin{probans}
$\beta$が存在するとせよ. $\alpha=:\alpha_1$の共役を$\alpha_2,\alpha_3\in \mb{R}$とする.$\mb{Q}(\alpha_1,\alpha_2,\alpha_3) = \mb{Q}(\beta,\alpha_2,\alpha_3)$であり, 左辺は$\mb{Q}$の正規拡大でかつ$\mb{R}$に含まれる体である.よって右辺も正規拡大かつ$\mb{R}$に含まれるので, とくに$\beta$の共役元はすべて実数である.しかし, $\beta^{3}\in \mb{Q}$から$\beta$の共役元には虚数も含まれるので矛盾である.\qed
\end{probans}


}



\clearpage
\section{有理数体上の拡大体}
最も出る問題は有理数体上の問題である.
\subsection{雑題}

神戸大のガロア群の計算問題が標準的である.
\begin{egprob} (H30神戸大B) \label{egprob:H30_kobe_B}
$f(x) = x^6-8$に対し, その最小分解体を$F$とする.
\ben
\item $f(x)$を$\mb{C}$で一次式に分解せよ.
\item $[F:\mb{Q}]$を求めよ.
\item $\gal{F/\mb{Q}}$の構造を求めよ.
\item 拡大$F/\mb{Q}$の$F,\mb{Q}$以外の中間体をすべて求めよ.
\een
\end{egprob}
\begin{egans}
(1):
\[f(x) = (x-\sqrt{2})(x-\sqrt{2}\omega)(x-\sqrt{2}\omega^2)(x+\sqrt{2})(x+\sqrt{2}\omega)(x+\sqrt{2}\omega^2)\]
(2):$F=\mb{Q}(\sqrt{2},\sqrt{-3})$が容易に分かる.$[\mb{Q}(\sqrt{2}):\mb{Q}]=2$, $[F:\mb{Q}(\sqrt{2})]=2$より$[F:\mb{Q}] = 4$

(3):Galであることをまず示さなければならない.分離性は完全体だから自明.生成元の共役根も$F$に属すから正規であることも自明.よってGal. 推進定理を使うと$(\mb{Z}/2\mb{Z})^{2}$.

(4):ひとつめの$\mb{Z}/2\mb{Z}$が$\gal{\mb{Q}(\sqrt{2})/\mb{Q}}$に対応するものであるとする.ガロア群の中間体で非自明なものは3つあることは容易.$\lan (1,0)\ran$に対応する体は, $\mb{Q}(\sqrt{-3})$である.$\lan (0,1)\ran$では$\mb{Q}(\sqrt{2})$である.$\lan (1,1) \ran$では$\mb{Q}(\sqrt{-6})$である.これらで全て.
\end{egans}


\begin{egprob} (\tf{有名問題}, \cite{aoyukie}問題4.1.16(1))
$p$を素数, $f(x)\in \mb{Q}[x]$とする.$f(x)$は$p-2$個の実数解と$2$個の虚数解を持つとする.$K$を$f$の最小分解体とするとき, $\gal{K/\mb{Q}} \cong \maf{S}_{p}$であることを証明せよ.
\end{egprob}
\begin{egans}
条件より複素共役$\sigma$はガロア群の元である. これは位数2である. ガロア群の元は$p$個の根の置換を引き起こし, その置換によって最小分解体から自身への$K$上の自己同形は一意に決定されるため, 単射準同型$\Phi: \gal{(K/\mb{Q})} \to \maf{S}_{p}$が引き起こされる. これによりガロア群を対称群の部分群に同一視する
. 既約性から$[L:K]$は$p$の倍数であり, Cauchyの定理という群論的事実から$\gal{L/K}$が位数$p$の元($p$-cycle)を持つことが従う. 一般に対称群$\maf{S}_n$は互換と$n$-cycleにより生成できるので$\gal{K/\mb{Q}}$が実は$\maf{S}_{p}$に一致する.
\end{egans}

\begin{prac} (\cite{aoyukie} 問題4.1.16(2))
$f(x) = x^5 - 4x + 2$の$\mb{Q}$上のガロア群は$\maf{S}_{5}$と同型であることを証明せよ.
\end{prac}

\tbox{
\begin{egprob} \label{prob:H25-II-1}
体$K=\mb{Q}(\sqrt{N},\sqrt{i+1})$が$\mb{Q}$上のGalois拡大体となるような最小の正の整数$N$と, そのときのGalois群$\gal{K/\mb{Q}}$を求めよ.
\end{egprob}
}{title=H25数学II-1}
\begin{egans}
$N\neq 1$を示す. $\alpha = \sqrt{i+1}$は偏角を$\pi/8$とする. $\beta = \ovl{\alpha} = \sqrt{1-i}$とすると$\alpha + \beta = \sqrt{2+\sqrt{2}}/2$が(半角公式より)分かるので$\sqrt{2+\sqrt{2}}\in K$だが, これは4次の根である(Check!). $K=\mb{Q}(\alpha)$も4次拡大なので$K=\mb{Q}(\sqrt{2+\sqrt{2}})$となるが, $K$は$\mb{R}$に含まれないのでおかしい. \par
$N=2$のとき, $K=\mb{Q}(\alpha,\sqrt{2})$は$\mb{Q}$上8次Galois拡大である. なぜなら$K\supsetneq \mb{Q}(\sqrt{2+\sqrt{2}})$で8次以上であるし, 8次以下であることも明らかだし, $\alpha\beta = \sqrt{2}$だからである. ガロア群を$G$とする. $\sigma\in G$は$\sigma(\sqrt{2})\in \set{\sqrt{2},-\sqrt{2}}$, $\sigma(\alpha) = \set{\alpha,-\alpha,\beta,-\beta}$であり, この選択で$\sigma$は一意に決まる. また, 8次だからいずれの選択も実現する元$\sigma$がある. \par
$\sigma_i$は$\sqrt{2}$を固定する元で$\sigma_1=\mr{id}$, $\sigma_2(\alpha) = -\alpha$, $\sigma_3(\alpha) = \beta$, $\sigma_4(\alpha) = -\beta$とする. $\tau$は$\tau(\sqrt{2}) = -\sqrt{2}$, $\tau(\alpha) = -\alpha$とする. このとき$G$の任意の元は$\sigma_i \tau^k$の形で書ける(Check!). たとえば$\tau \sigma_3\tau = \sigma_4$によって非アーベルがわかるので$D_4$と当たりをつけるのが妥当. 実際, $D_4$の生成式で$x = \tau$, $y=\tau\sigma_3$に取ると関係式が満たされる(Chekc!). よって$G\cong D_4$. \qed

\end{egans}
\begin{prac}
上の証明の細部をいくつか埋めよ.
\end{prac}

次も過去問だが, 比較的易しい.
\tbox{
\begin{egprob}
$x^4+9$の$\mb{Q}$上の最小分解体を$K$とするとき, $\mr{Gal}(K/\mb{Q})$を求めよ.
\end{egprob}
}{title=H19数学II}
\begin{egans}
最小分解体が$K=\mb{Q}(\sqrt{6}, i)$であることを示す. まず$x^4 + 9$の根は$\sqrt{3}\zeta_8^{t}$ ($t=1,3,5,7$)であり, $\mb{Q}(\sqrt{3}\zeta_8)\ni \sqrt{3}\zeta_8^{t}$なので
\bealn{
K&= \mb{Q}(\sqrt{3}\zeta_8) = \mb{Q}(\sqrt{3}\zeta_8, \sqrt{3}\zeta_8 + \sqrt{3}\zeta_8^7) \\
&= \mb{Q}(\sqrt{3}\zeta_8, \sqrt{6}) = \mb{Q}(\dfrac{1}{\sqrt{2}}\zeta_8, \sqrt{6}) \\
&= \mb{Q}(1+i,\sqrt{6}) = \mb{Q}(i,\sqrt{6})
}
よって$[K:\mb{Q}] = 4$である. ガロア群は位数4である. $\sigma\in \mr{Gal}(L/K)$は$\sigma(i) \in \set{i,-i}$と$\sigma(\sqrt{6}) \in \set{\sqrt{6}, -\sqrt{6}}$で決定され, 位数4だからすべての選択が可能である. $i$と$\sqrt{6}$の符号が変わるだけなので, ガロア群の任意の二つの元が可換であることが示せる. よってAbel群で, いずれの元も位数2なので$\cyc{2}\times \cyc{2}$であることが分かる.\qed
\end{egans}




\tbox{
\begin{egprob} \label{prob:H22-II-1}
$K=\mb{Q}(\sqrt{-17-4\sqrt{17}})$. $K/\mb{Q}$がガロア拡大であることを示し, $\mr{Gal}(K/\mb{Q})$を求めよ.
\end{egprob}
}{title=H22数学II}
\begin{egans}
$\alpha = \sqrt{-17-4\sqrt{17}}$とする.
$(\alpha^2 + 17)^2 -16\cdot 17 = \alpha^4 + 34\alpha^2 + 17 $であるから$T^4+34T^2+17$は$\alpha$を根に持つ. これは17に関してEis多項式で既約より, $K$は4次拡大. $\alpha$の$\mb{Q}$上共役元は$-\alpha, \beta:=\sqrt{-17 + 4\sqrt{17}}, -\beta$である. $\alpha,\beta$は虚である. $\alpha,\beta$の偏角は$\pi/2$とする. このとき
$\alpha\beta = - \sqrt{17}$である. $\alpha^2 \in K\implies \sqrt{17}\in K$なので$\beta\in K$が従う. よって$K=\mb{Q}(\alpha)$は$\alpha$の共役元をすべて持ち, $K/\mb{Q}$は4次ガロア拡大体. ガロア群を$G$とする. $G$は位数4だからアーベル群になる. $\sigma\in G$は$\sigma(\alpha)\in\set{\alpha,-\alpha,\beta,-\beta}$で決定され, この4通りがいずれも可能である. $\sigma(\alpha) = \alpha$なら$\sigma=\mr{id}$. $\sigma(\alpha) = -\alpha$なら$\sigma(\alpha^2) = \alpha^2$により$\sigma$が$\sqrt{17}$を固定すｒことが分かり, $\alpha\beta=-\sqrt{17}$から$\sigma(\beta)=-\beta$が分かる. $\sigma(\alpha) = \beta$ならば$\sigma(\alpha^2) = \beta^2$により$\sqrt{17}$の符号は反転し, $\sigma(\beta) = -\alpha$を得る. 同様に$\sigma(\alpha) = -\beta$なら$\sqrt{-17}$が反転し$\sigma(\beta) = \alpha$が分かる. 最後二つの元は明らかに位数4であるから$G$は$\cyc{4}$に同型である.
\end{egans}

\begin{rem}
なぜか知らないが, H23数学IIでも同じ問題が出ている(どころか, 誘導まで付いている).
\end{rem}






\tbox{

}{title=}



\begin{egprob}
$L$を$\mb{Q}(\sqrt[3]{3})$の$\mb{Q}$上のガロア閉包とするとき, $L/\mb{Q}$の中間体をすべて求めよ.
\end{egprob}
\begin{egans}
ガロア群が$\maf{S}_{3}$であることはすでに見た(cf.「Galois拡大」の節). これの中間体は, 非自明なものを見ると, 巡回部分群しかなく, 互換で生成されるものか$A_3$である. $A_3$は指数2で, これは$\tau$, つまり$\tau(\sqrt[3]{3}) = \sqrt[3]{3}\omega$, $\tau(\omega) = \omega$の固定体であるから$\mb{Q}(\omega)$である. \par
それぞれの互換が何を固定するかを調べよう. %つづく
\end{egans}








\subsection{4次既約多項式のガロア群}

\tbox{
\begin{egprob} \label{prob:H6-I-3}
$f(X) = X^4 + aX^2 + b$を$\mb{Q}$上の既約多項式とする. また$\mb{Q}$上での$f(X)$のガロア群を$G$で表す.
\ben
\item $b$が$\mb{Q}$の平方元なら, $G\cong \cyc{2}\times \cyc{2}$であることを示せ.
\item $\mb{Q}$において, $b$は平方元ではないが, $b(a^2-4b)$が平方元である場合には, $G\cong \cyc{4}$であることを示せ.
\item もし$b$も$b(a^2-4b)$も$\mb{Q}$において平方元でなければ, $G$は位数$8$の群であることを示せ.
\een
\end{egprob}
}{title=H6数学I-3}
\begin{egans}
$b=c^2$とする. $f(X)$の根は
\[\alpha = \sqrt{-\dfrac{a}{2} +\sqrt{\dfrac{a^2}{4} - b}},\quad  \beta = \sqrt{-\dfrac{a}{2} - \sqrt{\dfrac{a^2}{4} - b}}\]
を一つ選んで, $\pm \alpha, \pm \beta$ですべてである. $\alpha\beta = \sqrt{b}$であるように $\alpha,\beta$を選ぶことができる. $f(X)$の$\mb{Q}$上最小分解体を$K$とすると, $K=\mb{Q}(\alpha,\beta)$である. \\
(1): $\alpha\beta = c$であるように$c$を取り変える. $\beta = c/\alpha$なので, $K=\mb{Q}(\alpha)$である. とくに$[K:\mb{Q}] = |G| = 4$. $\sigma\in G$は$\sigma(\alpha)\in \set{\pm \alpha, \pm \dfrac{c}{\alpha}}$の行き先で決まり, どの元も位数2であるから$G$は巡回群ではないAbel群である. よって$G\cong \cyc{2}\times \cyc{2}$である. \\
(2): $b(a^2-4b) = d^2$とする. まず$b,d\neq 0$に注意しよう. 実際, $b=0$や$a^2-4b = 0$では$f(X)$の既約性に反することが分かるからだ. $\sqrt{b}\sqrt{a^2-4b} = d$であるように$d$を選ぶ. このとき,
\[\alpha^2 -\beta^2 = \dfrac{\sqrt{b}\sqrt{(a^2-4b)}}{\sqrt{b}} = \dfrac{d}{\sqrt{b}}\]
に注意. $\alpha\beta = \sqrt{b}$だから$K=\mb{Q}(\alpha,\sqrt{b})$である. $2\alpha^2 + a = \sqrt{a^2-4b}$なので$\sqrt{a^2-4b} \in \mb{Q}(\alpha)$であり, $d=\sqrt{b}\sqrt{a^2-4b}\in \mb{Q}$なので$\sqrt{b}\in \mb{Q}(\alpha)$である. よって$[K:\mb{Q}] = 4$, $|G| = 4$で, $\sigma\in G$のうち$\sigma:\alpha\mapsto -\beta$なるものが取れる. $\alpha^2-\beta^2 = \frac{d}{\sqrt{b}}$に$\sigma$を作用させたとき, 右辺は0にならず, $\sigma(\alpha) = -\beta$だから$\sigma(\beta) \in \set{\alpha,-\alpha}$である. このとき, 左辺の符号が変化するので$\sigma(\sqrt{b}) = -\sqrt{b}$である. よって$\alpha\beta = \sqrt{b}$から$\sigma(\beta) = \alpha$が従う. よって$\sigma^2(\alpha) =\sigma(-\beta) = -\alpha$なので$\sigma$は位数4の元でガロア群は$\cyc{4}$に同型である. \qed \\
(3): $G$が位数4の倍数かつ8以下なのは明らかで, $|G|=8$でないとするなら$|G| = 4$である. $\sqrt{b} \in \mb{Q}(\alpha)$である. $e=b(a^2-4b)\in \mb{Q}$とおく. $\sqrt{e} = p+q\sqrt{b}$なる$p,q\in \mb{Q}$がもし有れば, $e=p^2 + q^2b + 2pq\sqrt{b}$で, $pq=0$でなければならない. $q=0$なら$e$が平方元になり矛盾. $p=0$なら$e=q^2b$で, $b\neq 0$なので$a^2-4b = 0$である. これでは$f(X)$が平方完成されてしまうのであったから, 矛盾である. よって$K=\mb{Q}(\sqrt{b},\sqrt{e})$である. $d=a^2-4b$とおき, $e=bd$なので$K=\mb{Q}(\sqrt{b},\sqrt{d})$である. $\alpha^2 = -\frac{a}{2} + \frac{1}{2}\sqrt{d}$である. $\alpha = p+q\sqrt{b}+r\sqrt{d} + s\sqrt{bd}$とおく.
\bealn{
pq + rds = 0\\
ps + qr = 0\\
pr + qsb = \dfrac{1}{2}\\
p^2 + q^2b + r^2d + s^2bd = -\frac{a}{2}
}
だから, もし$qrs\neq 0$なら, $x=\dfrac{p}{r}, y=\dfrac{q}{s}$とすると$y\neq 0$であって,
\[\dfrac{pq}{rs} = xy = -d,\quad \dfrac{ps}{qr} = \dfrac{x}{y} = -1\]
であるから, $y=-x$で, $d=x^2$となるが, $d=a^2-4b$は平方元になってはならないので矛盾である. よって$qrs = 0$である. \\
\fbox{$q=0$のとき} $pr = 1/2$, $ps = rds = 0$ により$s=0$だがこのとき$\alpha$は4次既約多項式の根ではないので矛盾. \\
\fbox{$r=0$のとき} $pq=ps=0$, $qsb = 1/2$より$p=0$で, $\alpha = \sqrt{b}(q+s\sqrt{d})$である. $\alpha^2 = b(q^2 + s^2d) + 2bqs\sqrt{d}$なので,
\[q^2 + s^2d = -\frac{a}{2b},\quad qs\sqrt{d} = \frac{\sqrt{d}}{4b}\]
であるから, $q^2$と$s^2d$は$T^2 + \frac{a}{2b}T + \frac{d}{16b^2}$の\tf{有理数解}である. 特にこれの判別式が有理数の平方元であるべきだが, 実際は$\frac{a^2 -d}{4b^2} = \frac{a^2-(a^2-4b)}{4b^2} = \frac{1}{b}$だから, 仮定より平方元になっておらず矛盾である. \\
\fbox{$s=0$のとき} $pq=0, qr=0, pr=1/2$により, $q=0$を得るがこのとき$\alpha=p+r\sqrt{d}$は既約4次多項式の根になり得ない. \par
以上より, $\alpha\notin \mb{Q}(\sqrt{b},\sqrt{d})$が示されたので矛盾である. よって$|G|=8$. \qed
\end{egans}



\subsection{円分体}
\begin{prop}
\ben
\item $\gal{\mb{Q}(\zeta_n)/\mb{Q}}\cong (\mb{Z}/n\mb{Z})^{\times}$
\een
\end{prop}

\begin{eg} 任意の自然数$n$に対し, 有理数体上のガロア拡大で$\mb{Z}/n\mb{Z}$に同型なガロア群を持つものが存在することを示す. 算術級数定理により$p-1$が$n$で割れるような素数$p$が(無限個)存在するので, それを一つ取る. このとき$\mb{Q}(\zeta_{p})/\mb{Q}$のガロア群が$\mb{Z}/(p-1)\mb{Z}$である. この群の指数$n$の巡回部分群$H$が存在し, 基本定理から$H$に対応する中間体$M$があるが, $M/\mb{Q}$は$H$が正規部分群なのでガロア拡大になる. そして, $\gal{M/\mb{Q}} \cong \gal{(\mb{Q}(\zeta_n)/\mb{Q})}/H \cong \mb{Z}/n\mb{Z}$となる.
\end{eg}




\begin{comment}
%\subsection{($\spadesuit$)整数環の基底と絶対判別式}
\begin{defn}
$K$を代数体とする.$K$は自然に$\mb{Z}$を含んでいるが, $K$の元で$\mb{Z}$上整なもの全体(つまり$K$における$\mb{Z}$の整閉包)を$K$の\tf{整数環}といい, $\mac{O}_{K}$で表す.$\mac{O}_{K}$の$\mb{Z}$基底を\tf{整基底}といい, 整基底$\{ \omega_{i}\}_{i=1,\cdots, n}$ ($[K:\mb{Q}]=n$)に対して$\mb{Q}$係数の行列
\[(\mr{Tr}(\omega_{i}\omega_j))_{i,j}\]
の行列式を絶対判別式といい, $\Delta_{K}$と書く.$\Delta_{K}$は代数体の不変量である.
\end{defn}

$p$進完備化によって, 代数体$K$の整数環$\mac{O}_{K}$の基底や絶対判別式$\Delta_{K}$を決定することができるようになる.
\end{comment}





\subsection{$\cos{\dfrac{2\pi}{n}}$を添加した体}
$2\cos{\dfrac{2\pi}{n}} = \zeta_n + \zeta_n^{-1}$であり, $\zeta_n$は
\[T^2 - (\zeta_n + \zeta_n^{-1})T + 1\]
の根であるから$[\mb{Q}(\zeta_n): \mb{Q}(\zeta_n + \zeta_n^{-1})]\leq 2$である. また, $\mb{Q}(\zeta_n)/\mb{Q}$のガロア群の元$(\zeta_n\mapsto \zeta_{n}^{-1})$によりこの値は固定される.
したがって, $\mb{Q}(\zeta_n)\supsetneq \mb{Q}(\zeta_n + \zeta_n^{-1})$ は二次拡大である. すなわち, 円分体の中間体として現れる体である. このような体は演習問題に頻繁に現れるので, これについて調べよう. \par
まず, $2\cos{\dfrac{2\pi}{n}}$の$\mb{Q}$上の最小多項式を調べる際に三角関数の公式を用いることがあるので, 高校数学だがこれをおさらいしておこう.

\begin{prop}
\ben
\item $2\cos{2\theta} = (2\cos{\theta})^2 - 2$
\item $2\cos{3\theta} = (2\cos{\theta})^{3} - 3(2\cos{\theta})$.
\item
\een
\end{prop}

\begin{egprob} (\cite{aoyukie} 問題4.7.4)
$t=\zeta_7 + \zeta_7^{-1}$が満たす$\mb{Q}$上の三次方程式を求めよ.
\end{egprob}
\begin{egans} $\theta = \dfrac{2\pi}{7}$とおく.
$\gal{(\mb{Q}(\zeta_7)/\mb{Q})}\cong (\mb{Z}/7\mb{Z})^{\times}$の生成元は$\zeta_7 \mapsto \zeta_7^{3}$であるから, これにより$t$を写し続けることで$t$の$\mb{Q}$上共役元は
\[2\cos{2\theta} = \zeta_7^{2} + \zeta_7^{-2} = t^2-2,\quad 2\cos{3\theta} = \zeta_{7}^{3} + \zeta_7^{-3} = t^3 - 3t\]
となる. $1+\zeta_7 + \dots + \zeta_7^{6} = 0$なので
\[1 + t + (t^2 - 2) + (t^3-3t) = 0\]
これより$t^3 + t^2 -2t -1 = 0$.
\end{egans}
\begin{egans} (相反方程式的解法) 三角関数がだるかったら次のようにしてもよい.  $1+\zeta_7 + \dots + \zeta_7^6 = 0$なので, これを$\zeta_7^3$で割ると
\[(\zeta_7^3 + \zeta_7^{-3}) + (\zeta_7 + \zeta_7^{-2}) + (\zeta_7 + \zeta_7^{-1}) + 1 = 0\]
$t$の多項式で$\zeta_7^{k} + \zeta_7^{-k}$を記述できるので, その計算をすれば求められる.
\end{egans}

\begin{egprob}
$\alpha=2\cos{\dfrac{\pi}{9}}$の$\mb{Q}$上最小多項式を求めよ.
\end{egprob}
\begin{egans} (参考:2020代数II, 10/20練習問題)
3倍角公式のほうで\bolm{\theta = \dfrac{\pi}{9}}と置いて, $\cos{\dfrac{\pi}{3}} = \dfrac{1}{2}$なので, ただちに$T^3- 3T-1 = 0$を得る. これの最小性を示す方法は二つある. まず, $T=1$を代入すると$-3$であるから,  $T$を$T+1$に置き換えればEis多項式になると予想され, 実際$T^3 + 3T^2 - 3$なのでそうである. あるいは, これを$\mb{Z}$上で見たときに整数根を持たないことを論じるか, $\mb{F}_{2}$上に還元して$T^3 + T^2 + 1$が零点を持たないことを論じるかにより, これが$\mb{Z}[T]$で既約であることが分かるから, Gaussの補題により$\mb{Q}[T]$でも既約と言える.
\end{egans}


\begin{egprob} (2020代数学II, 12/15練習問題)
$\sqrt{5}\in \mb{Q}(\zeta_5)$を示せ. また, $\sqrt{-5}\notin \mb{Q}(\zeta_5)$を示せ.
\end{egprob}
\begin{egans} ガロア群が$\mb{Z}/4\mb{Z}$なので二次拡大は一つである. よって, その二次拡大を求めればよいと予想される.
$t=\zeta_5 + \zeta_5^{-1}$とするとこれは$\mr{id}$と$\zeta_5\mapsto \zeta_{5}^{-1}$で固定される元なので$\mb{Q}(t)$がまさしく2次拡大体である. $t$の$\mb{Q}$上の共役元は, $\zeta_5\mapsto \zeta_5^2$で移した$s=\zeta_5^2 + \zeta_5^{-2}$であるから, $(X-t)(X-s)$を考えると
\[t+s = -1, \quad ts = -1\]
が計算で得られ, $(X-t)(X-s) = X^2 + X - 1$である. よって$t,s$は$X^2 + X - 1$の根なので$\dfrac{-1\pm \sqrt{5}}{2}$であり, $\mb{Q}(t) = \mb{Q}(\sqrt{5})$を得る. これと$\mb{Q}(\sqrt{-5})$は異なるので$\sqrt{-5}$は属さない. \qed
\end{egans}


\begin{prac}前の問題の「相反方程式的解法」を真似して上の問題をもう一度解いてみよう.
\end{prac}

\tbox{
\begin{egprob} (2020代数学II, 12/22練習問題)
$n$を奇数とする. $\gamma:=\sin{\dfrac{2\pi}{n}}\notin \mb{Q}(\cos{\dfrac{2\pi}{n}})$を示せ.
\end{egprob}
}{}
\begin{egans}次に注意せよ!
\[\bolm{\gamma = \cos{\parena{\dfrac{\pi}{2} - \dfrac{2\pi}{n}}} = \cos{\dfrac{2(n-4)\pi}{4n}}}\]
$n$は奇数なので$n-4$と$4n$は互いに素である. よって, 右辺は$\mb{Q}(\zeta_{4n})$において$\cos{\dfrac{2\pi}{4n}}$と$\mb{Q}$上共役である. よって,
\[ [\mb{Q}(\gamma):\mb{Q}] = [\mb{Q}(\cos{\dfrac{2\pi}{4n}}):\mb{Q}] = \dfrac{\phi(4n)}{2}  \]
である. $\phi(4n) = 2\phi(n)$なので
\[[\mb{Q}(\gamma):\mb{Q}] = \phi(n) > \dfrac{\phi(n)}{2} = [\mb{Q}(\cos{\dfrac{2\pi}{n}}): \mb{Q}] \]
により$\gamma \notin \mb{Q}(\cos{\dfrac{2\pi}{n}})$でなければならない. \qed
\end{egans}

\begin{egprob} (\cite{aoyukie} 問題4.7.5, 改) \label{egprob:gal=z/3z_example}
$\gal{(\mb{Q}(\cos{\dfrac{2\pi}{9}})/\mb{Q})}$は何か?
\end{egprob}
\begin{point}
円分体に埋め込めば, アーベル拡大の中間体として取り扱える.
\end{point}
\begin{egans}
$t=2\cos{\dfrac{2\pi}{9}} = \zeta_9 + \zeta_9^{-1}$とおく. $t$は$t^3 - 3t - 1= 0$を満たす(略). よって$[\mb{Q}(t):\mb{Q}]\leq 3$である. 一方, $\zeta_9$が$T - tT + 1 = 0$の根なので$[\mb{Q}(\zeta_9):\mb{Q}]\leq 2$であり, これと$[\mb{Q}(\zeta_9): \mb{Q}] = 6$により, 不等式的に$[\mb{Q}(t):\mb{Q}] = 3$でなければならないことが分かる. 中間体$\mb{Q}(t)$に対応する部分群$H$の指数は3, つまり位数は2で, ガロア群の元$\sigma: \zeta_9\mapsto \zeta_9^{8}$が$\mb{Q}(t)$を固定することを考えると$H = \set{1,\sigma}$である. これを$\mb{Z}/6\mb{Z}$に翻訳すると, $H\cong \set{0,3}\subset \mb{Z}/6\mb{Z}$である. よって, ガロア対応から
\[\gal{(\mb{Q}(t)/\mb{Q})} \cong (\mb{Z}/6\mb{Z}) / H \cong \mb{Z}/3\mb{Z}.\]

\end{egans}


\subsection{演習問題}


\subsubsection{Level A}



\subsubsection{Level B}
\begin{prob} (H31専門科目3)
多項式$f(X) = X^4 + 6X^2 + 2\in \mb{Q}[X]$の$\mb{Q}$上の最小分解体を$K$とおく. $K$を$\mb{C}$の部分体とみなし,$F=K\cap \mb{R}$とおく.
\ben
\item $[K:\mb{Q}]$を求めよ.
\item $F/\mb{Q}$はガロア拡大であることを示せ.
\een
\end{prob}


\begin{prob} (H12数学II-1)
$\zeta = \exp{2\pi i/85}\in \mb{C}$, $K=\mb{Q}(\zeta)\cap \mb{R}$とするとき, $K$の$\mb{Q}$上のガロア群$\gal{K/\mb{Q}}$を求めよ.
\end{prob}

\begin{prob} (H2数学II-1)
$K=\mb{Q}(\sqrt[6]{3} + \sqrt{-3})$について次の問に答えよ.
\ben
\item $\mb{Q}$上の$K$のガロア閉包$\witi{K}$は$K$と一致するか?
\item $\gal{(\witi{K}/\mb{Q})}$を求めよ.
\item $\witi{K}$の部分体のうち$\mb{Q}$上6次拡大のものは?
\een
\end{prob}


\subsubsection{Level C}

\subsubsection{Hint・Answer}
{\small
\begin{probans}
(1): $[K:\mb{Q}] = 8$\\
(2): $F=\mb{Q}(\sqrt{3},\sqrt{7})$を拡大次数の比較から示せばよい.
\end{probans}

\begin{probans}
$K\supset \mb{Q}(\cos{\frac{2\pi }{85}})$で, $[\mb{Q}(\zeta): \mb{Q}(\cos{\frac{2\pi i}{85}})] = 2$かつ$K\neq \mb{Q}(\zeta)$なので$K=\mb{Q}(\cos{\frac{2\pi }{85}})$である. 円分体はアーベル拡大で, $K$に対応する部分群$H_K = \gal{(\mb{Q}(\zeta))/K}$は必ず正規. このとき$\gal{(K/\mb{Q})}\cong \gal{(\mb{Q}(\zeta))}/H_K$である.  ガロア群は$(\cyc{85})^{\times}\cong (\cyc{5})^{\times} \times (\cyc{17})^{\times} \cong \cyc{4}\times \cyc{16}$であり, $\zeta\mapsto \zeta^a$なる自己同型と$a\in (\cyc{85})^{\times}$が対応する. $\zeta + \zeta^{-1}$を固定する元が$a=1,-1$のものに限るから, $H_K = \set{1,-1}\subset (\cyc{85})^{\times}$である. これは$(\cyc{5})^{\times}\times (\cyc{17})^{\times}$において$\set{(1,1), (-1,-1)}$に対応し, 生成元は$(2,0)$と$(0,6)$なので, $2^2 \equiv -1\pmod{5}$と$6^8 \equiv -1\pmod{17}$に注意すると, $H_K$は$\cyc{4}\times \cyc{16}$の$\set{(0,0),(2,8)} = 2\mb{Z}/4\mb{Z}\times 8\mb{Z}/16\mb{Z}$に対応する. したがって
\[\gal{(K/\mb{Q})} \cong \mb{Z}/2\mb{Z}\times \mb{Z}/8\mb{Z}\]
\end{probans}


\begin{probans}
(2): $\witi{K} = \mb{Q}(\sqrt[6]{3}, \sqrt{-3})$は明らかであろう. これは位数12である. $\sigma\in \gal{(\witi{K}/\mb{Q})}$は$\sigma(\sqrt{-3}) \in \set{\pm \sqrt{-3}}$と$\sigma(\sqrt[6]{3})\in \set{\zeta^k \sqrt[6]{3}}$, $\zeta = \frac{1+\sqrt{-3}}{2}$で決まる. $\sigma: (\sqrt[6]{3},\sqrt{-3}) = (\zeta\sqrt[6]{3}, \sqrt{-3})$, $\tau:(\sqrt[6]{3},\sqrt{-3}) \to (\sqrt[6]{3}, -\sqrt{-3})$とすれば$\tau$は位数2, $\sigma$は($\zeta$を不変にするので)位数6であり, $\tau\sigma\tau = \sigma^{-1}$が分かるから $D_{6}$に同型. \\
(1): $\gal{\witi{K}/\mb{Q}}$のすべての元で$\alpha=\sqrt[6]{3} + \sqrt{-3}$が固定されたなら$\witi{K} = K$である. 実際は7個の元で固定されることを見ればよく, それはすぐ分かる. よって$K=\witi{K}$.
\\
(3): $D_6$の位数2の元の個数は$7$である. それは$\sigma^3$, $\sigma^k \tau$ ($k=0,1,\dots, 6$)である. あとは作業.
\end{probans}


}

\clearpage

\section{有限体}




有限体について覚えておくべきことをまとめる.$\mb{F}_{p^{n}}$を$x^{p^{n}} - x$の$\mb{F}_{p}$上の最小分解体と定義する.
\begin{prop}
$K$は位数が有限の体とする.代数閉包$\Omega = \ovl{\mb{F}_{p}}$を一つ固定する.以下が成り立つ.
\ben
\item $K$の位数は素数の冪である.
\item $L = \{ \alpha\in \Omega\mid \alpha^{p^{n}} = \alpha \}$は$x^{q}-x$の最小分解体である(したがって$\mb{F}_{p^{n}}$に同型である).
\item $K$はある$\mb{F}_{p^{n}}$に同型である.
\item $\mr{Frob}_{p^{n}}:\mb{F}_{p^{n}} \to \mb{F}_{p^{n}}$を$\mr{Frob}_{p^{n}}(x) = x^{p}$と定めると$\mr{Frob}_{p^{n}}$は体準同型である.これをFrobenius準同型という.
\item $\mb{F}_{p^{n}}/\mb{F}_{p}$はガロア拡大であり, $\gal{\mb{F}_{p^{n}}/\mb{F}_{p}}\cong \mb{Z}/n\mb{Z}$である.そしてその生成元は$\mr{Frob}_{p^{n}}$である.
\item $\mb{F}_{p^{n}}^{\times}$は位数$p^{n} - 1$の巡回群である.
\item 有限体は完全体である.
\een
\end{prop}

\prf (1):$K$の標数は素数$p$であり, これによって$K\supset \mb{F}_{p}$とみなせ, $[K:\mb{F}_{p}] = n$とすると$|K| = p^{n}$となる.\qed

(2):$L$が体であることは容易.$L$は$\Omega$の$\mb{F}_{p}$を含む部分体である.$\mb{F}_{p}\subset F\subset \Omega$が$x^{q} - x$の最小分解体であるなら, $F$には$x^{q}-x$の根が含まれているので, $L\subset F$である.$L$自身体であるから, $L=F$でなければならない.

(3):$K$の位数が$q=p^{n}$であるとする.$\mb{F}_{p}\subset K \subset \ovl{\mb{F}_{p}}$である.$L$と$K$の位数はいま$q$で等しい.$\mb{F}_{p}\subset L\subset \ovl{\mb{F}_{p}}$でもあるので, 任意の$x\in K$に対して$x^{q} = x$であることを示せば$K\subset L$かつ$|K| = |L|$が従うので$K=L$が得られる.$x^{q} = x$を満たすことについてだが, $K^{\times}$が位数$q-1$の群であることからLagrangeの定理によって$x^{q-1} - 1 = 0$なので直ちに得られる.\qed



\subsection{有限体上の既約多項式}
\begin{egprob}
$\mb{F}_{3}$上の$2$次モニック既約多項式をすべて求めよ.
\end{egprob}
\begin{screen}
意外にもこれには列挙以外に効率的な方法がある！この方法なら, 少し見つけてくるだけでよい.\tf{代数的な根がある$\mb{F}_{p^{n}}$に含まれてしまっているという統制のせいで, $\theta^{p^{n}} - \theta=0$という関係式が必ずあるのがポイントである.}
\end{screen}
\begin{egans}
$f(x)$を二次モニック既約多項式として, $f(x)$の根$\theta$を取る.$\mb{F}_{3}(\theta) = \mb{F}_{9}$なので, $3^{2} = 9$より, $\theta^{9}-\theta = 0$である.よって$f(x)$は$x^{9}-x$の因数である.$x^{9}-x$は$x(x-1)(x+1)$で割り切れる.
\[\dfrac{x^{9} - x}{x(x-1)(x+1)} = x^{6} + x^{4} + x^{2} + 1\]
$f(x)$はこの右辺の因数になっているような2次多項式というわけである.\par
さて, $x^2+1$が既約2次であることは簡単である.
\[\dfrac{x^{6}+x^{4} + x^2+1}{x^2+1} = x^4 + 1\]
$x^2-x-1$も既約2次なので
\[x^{4} + 1 = (x^2-x-1)(x^2+x+2)\]
以上より
\[\bolm{x^2+1,\quad x^2-x-1,\quad x^2+x^1}\]
\end{egans}




\subsection{($\spadesuit$)純非分離拡大}






\subsection{例題集}





\begin{egprob}
$f\in \mb{Z}[x]$の係数を$\bmod{p}$したものが$\mb{F}_{p}$上既約であるとき, $f$自身も$\mb{Q}$上既約であることが知られている(\cite{aoyukie}, 1.11節).一方でこの逆は成り立たない.$x^{4}+1$は既約であることが知られている.一方で,これを$\mb{F}_{p}$の多項式とみなしたものは常に可約であることを示す.

$p=2$のときは$x^{4}+1=(x+1)^4$なのでよい.$p>2$とする.$p^2-1$は$8$の倍数なので$x^{p^{2} - 1} - 1$は$x^{8} - 1$で割り切れ, $x^{4}+1$で割り切れる.よって$x^{4}+1$の根はすべて$x^{p^2} - x$の根であるから, $x^{4}+1$の任意の根$\alpha$について拡大次数$[\mb{F}_{p}(\alpha):\mb{F}_{p}] \leq 2$が従う.これは$x^{4}+1$が可約であることを示している.\bbox
\end{egprob}

\begin{egprob} (京大H20数学I)\\
$\ell$を素数とする.$\mb{F}_{p}$係数の$\ell$次既約モニック多項式の個数を求めよ.
\end{egprob}




\subsection{演習問題Part1} %雑多に
\subsubsection{Level A}
\begin{prob}
$\mb{Q}$代数として$\mb{Q}(\sqrt[6]{2})$と$\mb{Q}(\sqrt[6]{2}\omega , \sqrt[6]{2}\ovl{\omega})$は同型ではないことを示せ.ただし, $\omega=(-1+\sqrt{-3})/2$.
\end{prob}

\subsubsection{Level B}

\subsubsection{Level C}
\begin{prob} (\cite{neukirch}練習問題, Stickelbergerの定理)\\
代数体$K$について, その判別式を$d_K$とするとき, $d_K\equiv 0,1\pmod{4}$であることを証明せよ.
\end{prob}


\subsubsection{HINTS}
{\small
\begin{hint}
$\mb{Q}$代数の定義は覚えているか? 体論の問題というより環論の問題である.
\end{hint}

}

\newpage
\section{有理関数体}
\subsection{基本対称式で生成される体}
Artinの定理の応用例としてあげたように, 次が成り立つ.
\begin{prop} $k$を体とし, $n$変数有理関数体$L=k(x_1,\dots, x_n)$に変数の置換で$G=\maf{S}_n$を作用させる. このとき, $L^{G} = k(s_1,\dots, s_n)$. ただし, $s_i$は$i$次の$x_1,\dots, x_n$に関する基本対称式.
\end{prop}


\subsection{計算問題}
京大では少ないが, 東大ではよく有理関数体の拡大を見るので, それについての問題を挙げていく.

\lefba{
\begin{egprob} (マシュマロより)
$K=\mb{F}_{3}(T)$, $f(X) = X^3 - X\in K[X]$とする. $f(f(X))-T$の$K$上の最小分解体を$L$とするとき, $[L:K]$を求めよ. また, $L/K$の中間体で$K$上9次のものはいくつあるか.
\end{egprob}
}
\begin{egans}
$f(f(X)) - T = X^9-X^3 - (X^3-X) - T = X^9 - 2X^3 + X - T$である. $X^3 - X - T $の一つの根を$\alpha$とすると, $\alpha+1, \alpha-1$も根であるから
\[X^3 - X - T = (X-\alpha)(X-\alpha-1)(X-\alpha+1)\]
と因数分解される. よって$f(f(X)) - T$の根は$f(X) = \alpha, \alpha+1,\alpha-1$を満たす$X$である. 同様にして
\bealn{
X^3 - X - T - \alpha = (X-\beta)(X-\beta - 1)(X-\beta + 1)
}
と表す. $X^3 - X - 1$は$\mb{F}_{3}$上既約である. この根$r$を一つ取ると$r\in \mb{F}_{27}$である. $\mb{F}_{27}$の元$r$で$r^3 - r - 1 = 0$を満たすものについて   上の右辺で$X=\beta_0 + r$を代入すると$r(r-1)(r+1) = r^3 - r = 1$となるから, $f(\beta_0 + r) = \alpha + 1$であり, $f(f(\beta_0 + r)) - T=0$である. 同様に, $\mb{F}_{3}$上の既約多項式$X^3 - X + 1$でこれを考え, $\mb{F}_{27}$の元$s$で$s^3 - s + 1 = 0$を満たすものについて$f(f(\beta_0 + r)) - T = f(\alpha - 1) - T = 0$となる. \par
$X^3-X-1$の根$r$を取ると$r\pm 1$もまた根であり, $2r$, $2r\pm 1$は$X^3 - X + 1$の根なので, $f(f(X)) - T$の根は
\[\beta + k, \beta + r + k , \beta + 2r + k,\quad (k=-1,0,1)\]
である. したがって最小分解体$L$は
\[L = \mb{F}_{3}(T)(\beta, r) = \mb{F}_{27}(T)(\beta) = \mb{F}_{27}(T,\alpha)[X]/(X^3-X-T-\alpha)\]
である. 一旦$[L:K] = 27$であることを認める. $f(f(X)) - T$が分離多項式であることは微分をすれば分かる. よって$L/\mb{F}_3(T)$は分離拡大だからガロア拡大. $G = \gal{(L/\mb{F}_{3}(T))}$とする. $\sigma\in G$は$\sigma(\beta)\in \set{\beta + mr + k \mid m,k\in \mb{F}_{3}}$と$\sigma(r) \in \set{r,r+1,r-1}$で決まる. 求めるべき個数は, $G$の位数3の部分群の個数であるから, 位数3の元の半分の個数を求めればよい. 実は, $G$に位数9以上の元は存在しない. 実際, $\sigma\in G$が
\[\sigma(\beta) = \beta + m r + k,\quad \sigma(r) = r+a\]
だとすれば,
\[\sigma^3(r) = r+3a = r\]
\bealn{
\sigma^3(\beta) &= \sigma^2(\beta + mr + k) \\
&= \sigma((\beta + mr + k) + m(r+a) + k )\\
&= \sigma(\beta + 2mr + ma + 2k) \\
&= (\beta + mr + k) + 2m(r+a) + ma + 2k \\
&= \beta + 3mr + 3ma + 3k = \beta
}
であるから$\sigma$の位数は3以下である. したがって$G$の位数3の元は26個であるから, 9次中間体の個数は$13$個. \\
\fbox{$[L:K] = 27$を示す}

\end{egans}








\lefba{
\begin{egprob} (東大専門2021)
$K=\mb{C}(X,Y,Z,W)$に$G=\maf{S}_4$を変数の置換で作用させる. $F=K^G$とおき, $L=K(\sqrt{X})$の$F$上のガロア閉包を$M$とする.
\ben
\item $[M:K]$を求めよ. $M=K(x)$をみたす$x\in M$を求めよ.
\item $K$と$M$の中間体の個数を求めよ.
\item $K$と$M$の中間体$E\neq K,M$のうちで$F$のガロア拡大であるものをすべて求めよ. 各々の$E$に対して$E=K(y)$をみたす$y\in E$を一つ求めよ.
\een
\end{egprob}
}
\begin{egans}
{\small
(1): $s_i$ ($i=1,2,3,4$)を$X,Y,Z,W$に関する$i$次基本対称式とすると$F=\mb{C}(s_1,s_2,s_3,s_4)$である.
\[(T^2 - X)(T^2-Y)(T^2-Z)(T^2-W) \in F[T]\]
である. 実際, $T^8 - s_1T^6 + s_2T^4 -s_3 T^2 + s_4 $になるから. そしてこれは$F$上で既約である.
{\small
\lefba{ 少し自信がないが, 次のようにやったら既約が言えると思う. 上の多項式がもし可約であるなら, それは4次と4次の積にならねばならないことがわかる. なぜなら, 上の多項式を$\gal{(K/K^G)}\cong \maf{S}_4$の元で作用すると解の置換が起きるが, $X\mapsto Y$なる$\maf{S}_{4}$の元を作用させると$\sqrt{X}$は$\pm \sqrt{Y}$に移らなければならない. $Z,W$でも同様である. 符号については掴めないが, これで$\sqrt{X}$の$F$上の共役元が4つ以上あることがわかる. よって既約でないなら$\sqrt{X}\in \set{\pm \sqrt{X}, \dots, \pm\sqrt{W}}$の$\gal{K/K^{G}}$による作用の軌道が4つの元しか持たないことになり,  $(T- \sqrt{X})\dots (T\pm \sqrt{W})$の形の多項式が$\sqrt{X}$の$F$上最小多項式であるべきである. これは, 定数項が$\pm \sqrt{s_4}$なのでありえない. よって既約.
}
}
よって$M\supset K(\sqrt{X},\sqrt{Y},\sqrt{Z},\sqrt{W})$. 逆に右辺の体は$L$を含み, $\sqrt{X}$の$F$上の共役元をすべて含んでいるのでガロア拡大(標数0なので分離的). よって等号が成立. $[M:K] = 16$は明らかだろう. $M/F$がガロアで$F\subset K$なので$M/K$もガロア. $\gal{(M/K)}$は位数16の群であり, $\sigma\in \gal{M/K}$は$T^2-X,\dots, T^2-W\in K[T]$を変えないので, $\sigma$は4つの集合$\set{\sqrt{X},-\sqrt{X}},\dots, \set{\sqrt{W}, -\sqrt{W}}$に作用し, この作用の様子で$\sigma$は一意的に決まる. よって単射準同型$\gal{(M/K)} \to (\cyc{2})^4$を得る. 位数がともに16なので, 全射でもある. よって同型が引き起こされる. さて, $M$を$K$上の単拡大$K(x)$と見るには, すべての$\sigma\in \gal{(M/K)}$で$\sigma(x)\neq x$であるような$x\in M$を見つければよい(ガロア対応による). よって$x=\sqrt{X} + \sqrt{Y} + \sqrt{Z} + \sqrt{W}$とすれば16個の元すべてで$x$は固定されないことが確かめられる. \\
(2): $\mb{F}_{2}^4$の部分空間の個数を求めればよい. より一般に$\mb{F}_{p}^n$の$k$次元部分空間の個数について論じる. $k=0$なら1個である. $k\geq 1$とする.  任意の$k$個の独立なベクトルの順序対$(v_1,\dots, v_k)$は$k$次元部分空間$V$を張る. 異なる順序対に対して同じ$V$が得られることもあるが, そのいずれも$V\cong \mb{F}_{p}^{k}$の基底である. $\mb{F}_{p}^{k}$の順序を考慮した基底の個数は
\[N_k:=(p^k-1)(p^k-p)\dots (p^k-p^{k-1})\]
である. よって, $(v_1,\dots, v_k)\mapsto V$は$N_k$対1の写像であり, $k$個の独立なベクトルの順序対の個数は
$(p^n-1)(p^n-p)\dots (p^n-p^{k-1})$
である. 以上をまとめ, $k$次元部分空間$V$の個数は$n(p,k) := \prod_{i=0}^{k-1} \dfrac{p^n - p^{i}}{p^k - p^i}, n(p,0) = 1$
で得られる. これを$p=2$で考え,
\bealn{
&n(2,0) = n(2,4) = 1 \\
&n(2,1) = n(2,3) = 15\\
&n(2,2) = \dfrac{16-1}{4-1}\cdot \dfrac{16-2}{4-2} = 35
}
だから$1+1+15+15+35 =$ \bolm{67}個.\\
(3): $\sigma\in \gal{M/F}$の$K$への制限を$\phi\in \gal{K/F}$とする. $\phi$は$X,Y,Z,W$の置換を起こすが, $\phi(\sqrt{X}) = \sqrt{\phi(X)}$ などと定めることで$\phi$は$F$同型$M\to M$に拡張される.
このとき$\phi^{-1}\circ \sigma\in \gal{M/F}$は$X,Y,Z,W$を固定するので$\sigma\circ\phi^{-1}\in \gal{M/K}$であり, $\phi^{-1}\circ \sigma$は$\sqrt{X},\dots, \sqrt{W}$の符号のみを変化させる. $\alpha\in \gal{(M/K)}$であって, $\alpha(\sqrt{X}) = \epsilon_{X}\sqrt{X}$ ($\epsilon_X = \pm 1$)などとするとき, この$\epsilon_X,\dots, \epsilon_W\in \set{\pm 1}$によって$\alpha$は一意に決まるが,  この$\alpha$のことを$[\epsilon_X,\epsilon_Y,\epsilon_Z,\epsilon_W]$と表す. 先の考察で, $\sigma\in \gal{(M/F)}$が$\sigma = \phi\circ [\epsilon_X,\epsilon_Y,\epsilon_Z,\epsilon_W]$の形に表されることが分かった. \par
$\gal{(M/K)}$は$[\pm1, \pm1 , \pm1, \pm1]$ (複合任意)の全体の集合であり, $(\cyc{2})^4$に同型であったからアーベル群である.  $E$に対応する部分群を$H(E)=\gal{(M/E)}$とする. $M/E$がガロア拡大であることと$H(E)$が$\gal{M/F}$の正規部分群であることは同値であるから, 任意の$\sigma\in \gal{(M/F)}$に対して$\sigma H(E) \sigma^{-1} = H(E)$であるかを見る. $\sigma = \phi\circ [\epsilon_X,\dots, \epsilon_W]$と表すとき, $[\epsilon_X,\dots, \epsilon_W]H(E)[\epsilon_X,\dots,\epsilon_W]^{-1} = H(E)$なので($(\cyc{2})^4$がアーベル群だから), $\sigma=\phi$に限定して調べてもよい. すなわち, $\sigma$は$\sqrt{X},\dots, \sqrt{W}$をそれぞれ$\sqrt{\sigma(X)},\dots, \sqrt{\sigma(W)}$に移すものに限定してよい. さて, $[\epsilon_X,\dots \epsilon_W]\in \gal{M/K}$に対して
\[\sigma[\epsilon_X,\dots, \epsilon_W]\sigma^{-1} = [\epsilon_{\sigma^{-1}(X)},\dots, \epsilon_{\sigma^{-1}(W)}]\]
であることは容易に確かめられる. $\sigma$はいかなる文字の置換も起こせるので, ここから $\gal{(M/K)}$が$\gal{(M/F)}$のどのような共役類から成るかが分かる. すなわち, $\gal{(M/K)}$は次の5つの共役類の直和である.
\bealn{
&\set{\mr{id}},\quad \set{[-1,-1,-1,-1]},\\
C_1 &= \set{[a,b,c,d]\mid a,b,c,dのうち-1は一つ}\\
C_2 &= \set{[a,b,c,d]\mid a,b,c,dのうち-1は二つ}\\
C_3 &= \set{[a,b,c,d]\mid a,b,c,dのうち-1は三つ}
}
また, $|C_1| = |C_3| = 4, |C_2| = 6$である. $H(E)$は$\gal{(M/K)}$に含まれた$\gal{(M/F)}$の正規部分群だから, 上の5つの共役類のいくつかの合併でかつ群をなすものにならねばならない.  そのようなものは$\set{\mr{id}, [-1,-1,-1,-1]}$しかありえない(簡単に分かる). したがってこの部分群に対応する体が答え. その体$E$が$E':=K(\sqrt{XY},\sqrt{XZ},\sqrt{XW})$に一致することを示す. $E'$を固定する任意の$\sigma\in \gal{M/F}$を取れば, まずそれは$K$を固定するので, $\sigma$は変数の置換を起こさない. よって$\sigma=[\epsilon_X,\dots, \epsilon_W]$と表せるが, $\sigma$は$\sqrt{XY},\sqrt{XZ},\sqrt{XW}$も固定するので
\[\epsilon_X\epsilon_Y = \epsilon_X \epsilon_Z = \epsilon_X \epsilon_W = 1\]
を得る. 等式全体に$\epsilon_X$を掛けて$\epsilon_X = \epsilon_Y = \epsilon_Z = \epsilon_W$を得るので$\sigma\in H(E)$である. $H(E') \supset H(E)$は明らかなので, $H(E') = H(E)$が従い, $E = E'$である. 最後に, $y=\sqrt{XY} + \sqrt{XZ} + \sqrt{XW}\in E$に対して $E=K(y)$を示す. $H(E)\subset H(K(y))$は自明. 逆に$K(y)$を固定する$\sigma\in \gal(M/K)$を取る. $\sigma=[\epsilon_X,\dots, \epsilon_W]$とすると, $\sigma(y) = y$よりやはり$\epsilon_X\epsilon_Y = \epsilon_X\epsilon_Z = \epsilon_X\epsilon_W = 1$なので同様に$\sigma\in H(E)$となる.
ゆえに$E=K(y)$である. \qed

}
\end{egans}




\lefba{
\begin{egprob} (H13数学II) \label{prob:H13-II-1}
$K=\mb{C}(x_1,\dots, x_n)$の$\mb{C}$上の自己同型$\sigma$を
\[\sigma(x_i) = x_{i+1}\quad (1\leq i\leq n-1),\quad \sigma(x_n) = x_1\]
によって定義する.
\ben
\item $\sigma$による$K$の不変部分体$F=\set{\alpha\in K \mid \sigma(\alpha) = \alpha}$に対して, $K=F(\sqrt[n]{\alpha})$を満足する$\alpha\in F$をひとつ求めよ.
\item $n\geq 3$であれば, (1)の条件を満たす$\alpha\in F$は$x_1\dots, x_n$の対称式には取れないことを示せ.
\een
\end{egprob}
}
\begin{egans}
(1): $\lan \sigma \ran$ は$K/F$に忠実に作用する. アルティンの定理より$\gal{(K/F)}\cong \cyc{n}$である. $\alpha = (\sum_{j=1}^{n} \zeta_n^j x_j)^n$とおくとよいことを示す. $\sqrt[n]{\alpha} = \sum_{j=1}^{n} \zeta_n^j x_j$として$n$乗根を選ぶ. $\alpha\in F$は$\zeta_n^n = 1$なので分かる. $K_0 = F(\sqrt[n]{\alpha})$と$K$を比較しよう. $K_0$を固定する部分群が$\set{0}\subset \cyc{n}$であればよい. $\sigma^k(\sqrt[n]{\alpha}) = \zeta_n^{-k}\sqrt[n]{\alpha}$なので, 確かに固定部分群は$\set{0}$である. よって$K_0 = K$である. \qed \\
(2): $x_1,\dots, x_n$の第$i$次基本対称式$s_i$を考える. $L/M$には $\maf{S}_n$が変数の置換で作用しており, $L^{\maf{S}_n} = M$である. アルティンの定理から$L/M$はガロア拡大で$\gal{(L/M)}\cong \maf{S}_n$である. $n\geq 3$により$\lan \sigma \ran \subsetneq \maf{S}_n$であるから$M\subsetneq F$である. $[K:F] = n$なので, $[F:M] = (n-1)!$となる. $M(\sqrt[n]{\alpha})/M$は, 任意の$\sigma\in \maf{S}_n$で$\sigma(M(\sqrt[n]{\alpha})) = M(\sqrt[n]{\alpha})$なのでガロア拡大である. $K=F(\sqrt[n]{\alpha})$なので$\sqrt[n]{\alpha}$の$F$上最小多項式は$n$次であり, $\sqrt[n]{\alpha}$の$M$上最小多項式は$n$次以上である. $T^n - \alpha \in M[T]$なので, ちょうど$n$次でありこれが最小である. よって$[M(\sqrt[n]{\alpha}):M] = n$が分かる. さて, \ao{$M(\sqrt[n]{\alpha})/M$は$n$次ガロア拡大だからこれは$\maf{S}_n$の指数$n$の正規部分群に対応する.} $n\geq 3$かつ$n\neq 4$において, $\maf{S}_n$の自明でない正規部分群は$A_n$に限るが, これは指数が$2$なので$n=2$となり矛盾である. $n=4$においては, $\maf{S}_{4}$の正規部分群は$A_4$と$K$(Kleinの四元群) であり, これらは指数$2,6$であるから指数$4$になり得ない. よって矛盾である.

\end{egans}


\subsection{演習問題}
\subsubsection{Level B}


\begin{prob} (H20数学II-2)
$K=\mb{C}(t)$, $p,q\in \mb{C}$, $L=\mb{C}(t^3+pt+q)$とする. $K/L$は代数拡大であることを示し, さらに$K/L$のガロア閉包とその$L$上の拡大次数を求めよ.
\end{prob}



\begin{prob} (H5数学II-2)
$n\geq 2$, $L=\mb{C}(X)$, $K=\mb{C}(X^n + X^{-n})$とする.
\ben
\item $L$は$K$上のガロア拡大になることを示し, そのガロア群を求めよ.
\item $n=3$のとき$L/K$の中間体をすべて求めよ.
\een
\end{prob}



\subsubsection{Hint・Answer}
\begin{probans}
ガロア閉包を$\witi{K}$とする. $[\mb{C}(t):\mb{C}(t^3+pt+q)] = 3$は$T^3 - pT - q$の$\mb{C}(t)$上既約性から従う(仮に可約であれば一次の根を持つ. その根は整閉性から$\mb{C}[t]$に属すはずだが次数を見るとそれはあり得ない).
$p=0$のとき$\witi{K} = \mb{C}(t)$となるので$[\witi{K}:L] = 3$. $p\neq 0$なら6である.
\end{probans}

\begin{probans} $X$は
$T^{2n} -(X^n + X^{-n})T^n + 1$の根なので代数拡大. これの根は$\zeta^k X$と$\zeta^k X^{-1}$で, 正規拡大である. 標数0なので分離的. よってガロア拡大. $M=\mb{C}(X^n)$とする. $L\supset M \supset K$という列を考えると$[L:M] = n$である. $X\to X^{-1}$という自己同型$\mb{C}(X)\to \mb{C}(X)$で$K$の元はすべて固定されるが$M$はそうではないので$[M:K] = 2$である. $G = \gal{(L/K)}$とすると, $L=K(X)$であるから, $\sigma\in G$は$\sigma(X)\in \set{\zeta^k X, \zeta^k X^{-1}}_{k=0,1,\dots, n}$で決まる. \tf{これはアーベル群ではない.} $\sigma:X\mapsto \zeta X$, $\tau: X\mapsto X^{-1}$とする.
\[\tau \sigma \tau^{-1} = (X\mapsto \zeta^{-1}X) = \sigma^{-1}\]
なので, $G\cong D_{n} = \lan \sigma,\tau\mid \sigma^n = \tau^2 = e, \tau\sigma\tau^{-1} = \sigma^{-1} \ran$である. \\
(2): $D_3 = S_3$で部分群は6個. 2次拡大は$\mb{C}(X^n)$. 3次拡大は, $\mb{C}(X+X^{-1}), \mb{C}(X+\zeta_3 X^{-1}), \mb{C}(X+\zeta_3^{2}X^{-1})$.
\end{probans}




\clearpage
\section{($\heartsuit$)ガロア群}
ガロア群を計算する問題を取り扱う.
\subsection{群論パズル}
ガロア群を群論レベルの話に持ち込んで非自明な議論を行うことがある.このような部分で苦労するくらいなら, いっそ群論を極めたいものなのだが.ここでは, ガロア群を考えるときに重要になる群論レベルの話題についてまとめる.
\begin{prop} (ガロア群と群論まとめ)
以下, $K/\mb{Q}$は体拡大とする.
\ben
\item (基本) $S_n$は互換全体で生成できる.
\item (基本) $S_n$は互換一つと$n$-cycleで生成できる.
\item ($\spadesuit$) $A_p$は$3$次巡回置換一つと$p$-cycleで生成できる.
\item ($\spadesuit$) $S_{p}$は互換と$p$-cycleで生成できる.
\item $K$が$n$次の$f(x)\in \mb{Q}[x]$の最小分解体であるとき, 単射準同型$\gal(K/\mb{Q})\hookrightarrow S_{n}$が存在する.
\item $q$が素数で$\sigma\in S_n$が位数$q > \frac{n}{2}$のとき, $\sigma$は$q$次巡回置換である.
\een
\end{prop}
\prf
\ben
\item
\item
\item
\item (\href{https://math.stackexchange.com/questions/2792165/let-a-be-a-p-cycle-in-s-p-and-let-b-be-a-transposition-in-s-p-show-s/2792188}{Stack}) $c= (0123\cdots p-1)$, $t=(0,a)$, $a\in \mb{F}_{p}^{\times}$とする.$G$を$c,t$で生成される部分群とするとき, $ctc^{-1} = (1,a+1)$であり, $c^{k} t c^{-k} = (k,k+a)$となる.よって$(0,a),(a,2a),\cdots$はすべて$G$に属すが, $ab = 1$となる自然数$b$が存在し, すべての$k$で$(0,k)\in G$である(ここがポイント).$(i,j)\in G \iff (0,i)(i,j)(0,i)=(0,j)\in G$なので$(i,j)\in G$でありすべての互換を$G$が含むので$G=S_p$\qed
\item $j$が$\sigma$で固定されない元であるとするとき, $q$が素数だから, $j$は$\sigma$で$q$回作用して初めて$j$に戻る.このときの軌道が$q$次巡回置換となる.$k$がもうひとつの$\sigma$で固定されない元であるとき, もしこの軌道に属さなければ, $k$の$\lan \sigma \ran$による軌道も$q$次巡回置換を与え, ふたつの交わらない$q$次巡回置換を得る.$2q>n$であるからこの軌道は交わるが, これは$k$が最初の軌道に属さないことに反する(軌道分解は同値類分解である).
\een

\begin{eg} (H18数学II-1) \label{eg:H18_II_1}
$S_5$の部分群$H$の位数が$15$で割れるとする.このとき$H=A_5$または$S_5$であることを示す.\par
Cauchyの定理によって$H$は位数$3$の元と位数$5$の元を含むので, (5)より$H$は3-cycleと5-cycleを含む.よって(3)より$H\supset A_{5}$である.$H$の位数は$60$か$120$なので, $H=A_5$または$H=S_5$である.\par
これを用いてH18数学II-1を解く.$f(x)$の根の一つを$\alpha$とするとき$f$の既約性から$[\mb{Q}(\alpha):\mb{Q}] = 5$なので$[K:\mb{Q}]$は$5$で割れる.$f(X)$の$F$上の既約因子のうち$3$次のほうの根の一つを$\beta$とおくと$[F(\beta):F] = 3$なので, $G$の位数$=[K:\mb{Q}]$は$15$で割れる.$G=\gal(K/\mb{Q})$は$S_5$の部分群と同一視でき, $15$で割れるから$G \cong A_{5}$または$G\cong S_5$.\qed
\end{eg}

ガロア群を具体的に決定する状況において, アーベル群のうちはまだいいが非アーベルのものが出ると大変なので, 生成式を覚えておいて当たりをつけよう. とくに, \sout{性格の悪い大学では}2面体群くらいのものが出る可能性がそこそこある.
\tbox{
\begin{point}
\bealn{
D_n &= \lan \sigma,\tau \mid \sigma^n = \tau^2 = 1, \tau \sigma \tau = \sigma^{-1} \ran
}
\end{point}
}{title=群の生成式}
この説明だけでは抽象的に感じられるかもしれませんが, 次の表示を用いると具体的に理解しやすくなります.
\begin{prop} (東北大の問題を解いたときのアイデア)
$\maf{S}_{4}$に$(1234)$と$(24)$で生成される部分群$G$の位数は8であり, $D_{4}$に同型である. また, $G$は次の条件$(\ast)$を満たすような写像$\phi:\set{1,2,3,4}\to \set{1,2,3,4}$の全体の集合でもある.

\lefba{
$(\ast)$ $\phi(\set{1,3}) = \set{1,3}$または$\phi(\set{1,3}) = \set{2,4}$.
}
\end{prop}



\begin{egprob} (H2数学I-2)
次の命題はそれぞれ正しいか.
\ben
\item $H$が群$G$の部分群でその指数$n$が$2<n<\infty$のとき, $G$の正規部分群$N$で, $N\neq G$, $N$の$G$での指数$\leq n!/2$を満たすものがある.
\item $K$が体, $n$が偶数のとき, $K$の$n$次分離拡大体$L$は, 少なくとも一つ$K$の二次拡大体を含む.
\een
\end{egprob}
\begin{egans}
(1): $G=A_5$は非自明な正規部分群を持たない. $H=\lan (123)\ran$は指数3の部分群だが, 指数3以下の正規部分群は持たない. \\

(2): $L/K$をガロア群が$S_5$であるガロア拡大とする. $A_5$に対応する中間体$M$を取る. $L/M$は60次ガロア拡大である. $\gal{(L/M)}\cong A_5$である. $A_5$は単純群なので, 非自明な正規部分群を持たない. とくに, 指数2の部分群を持たない(指数2は正規だから). 指数2の部分群は, 下の体から2次なものに対応するので, $L/M$は二次拡大体を持たない.
\end{egans}


\subsection{計算問題}




\begin{egprob} (東大B?)
有理関数体$\mb{R}(X)$を$K$とする. $F=\mb{R}\parena{X^4 - \dfrac{1}{X^4}}$とし, $K$の$F$上のGalois閉包を$L$とおく.
\ben
\item $[L:F]$を求めよ.
\item $\mr{Gal}(L/F)$は位数8の元を持つことを示せ.
\item $L/F$の4次中間体をすべて求めよ.
\een
\end{egprob}
\begin{egans} (\href{https://detail.chiebukuro.yahoo.co.jp/qa/question_detail/q14198600154}{知恵袋})

\end{egans}










\begin{egprob} (\cite{aoyukie}問題4.1.17)

\end{egprob}
\begin{egprob} (\cite{aoyukie}問題4.7.5)
\ben
\item $t = \zeta_7 + \zeta_7^{6}$が満たす$\mb{Q}$上の三次方程式を求めよ.
\item $[\mb{Q}(\zeta_9): \mb{Z}(t)]\leq 2$を示すことにより$[\mb{Q}(t):\mb{Q}] = 3$を示せ.
\item $\gal{(\mb{Q}(t)/\mb{Q})}$は何か?
\een
\end{egprob}









\subsection{演習問題} %Galois群計算
\subsubsection{Level A}
\begin{prob}
$K$は$x^6+3$の$\mb{Q}$上最小分解体とする. $\gal{(K/\mb{Q})}$を求めよ.
\end{prob}


\subsubsection{Level B}
\begin{prob} (H24数学II-1)
$1$以上の整数$a$に対して, $K=\mb{Q}(\sqrt{3+a\sqrt{5}})$が$\mb{Q}$のGalois拡大体となるようなものを求め, そのような$a$に対してGalosi群$\mr{Gal}(K/\mb{Q})$を求めよ.
\end{prob}

\begin{prob} (H29専門科目3) \label{prob:H29senmon-3}
\ben
\item $S_3$を1,2,3に関する対称群と見て$S_3\subset S_5$とみなす. $\sigma = (123)\in S_5$, $G = \lan \sigma \ran$, $\tau = (45)\in S_5$, $h=\lan \tau\ran$とおく.
$N_{S_5}(G) = \set{\eta\in S_5 \mid \eta G \eta^{-1} = G}$ とおく. このとき$N_{S_5}(G) = S_3\times H$であることを示せ.
\item $f(X)\in \mb{Q}[X]$は5次で$K\subset \mb{C}$を$f(X)$の$\mb{Q}$上の最小分解体とする. $K$は$[K:F] = 3$となる部分体$F$がただひとつ存在すると仮定する. このとき$f(X)$は$\mb{Q}$係数の3次の既約多項式で割り切れることを示せ.
\een
\end{prob}

\subsubsection{Level C}

\begin{prob} (東工大)
$K=\mb{Q}(\sqrt{1+\sqrt{2}}, i)$とする.
\ben
\item $K/\mb{Q}$はGalois拡大であることを示し, そのガロア群を求めよ.
\item $K/\mb{Q}$のすべての中間体を求めよ.
\een
\end{prob}


\begin{prob} (R2専門科目2) \label{prob:R2-senmon-2}
$p\geq 5$を素数とし, $\sqrt{-p}+\sqrt[3]{p}$を含む$\mb{Q}$のガロア拡大体のうち最小のものを$K$とする. このとき
$\mr{Gal}(K/\mb{Q})$を求めよ.
\end{prob}

\begin{prob} (H30専門科目3) \label{prob:H30senmon-3}
$X^7-11$の$\mb{Q}$上最小分解体を$K$とする.
\ben
\item $[K:\mb{Q}]$を求めよ.
\item $K/\mb{Q}$の$\mb{Q},K$でもない真の部分群の個数を求めよ.
\item (2)の部分群のうち正規部分群であるものの個数を求めよ.
\een

\end{prob}


\begin{prob} (\tf{代数学の基本定理})

\end{prob}


\newpage

\subsection{Hint・Answer}
{\small
\begin{probans}
$\alpha = \sqrt[6]{-3}$, $\zeta = \zeta_6$とする. $x^6 + 3$の根は$\alpha\zeta^k$なので$K=\mb{Q}(\alpha,\zeta)$. $2\zeta - 1 = \alpha^3$なので$K=\mb{Q}(\alpha)$. $x^6 + 3$はアイゼンシュタイン多項式なので$[K:\mb{Q}] = 6$. $\sigma\in \gal{(K/\mb{Q})}$は$\sigma(\alpha)\in \set{\alpha\zeta^k\mid k=0,1,\dots, 6}$の写し方で決まる. $\sigma(\alpha) = \alpha\zeta^k$ ($k=1,3,5$)ならば
\[\sigma(2\zeta) = \sigma(\alpha^3 + 1) = 1-\alpha^3 = 2\zeta^5\]
だから
\[\alpha\xmapsto{\sigma} \alpha\zeta^{k} \xmapsto{\sigma}\alpha\zeta^k \cdot \zeta^{5k} = \alpha \]
より$\sigma$の位数は2である. 位数2の元が3個見つかるのでガロア群は$\cyc{6}$ではあり得ず, $S_3$に同型である($k=1,3,5$のものが互換に対応).
\end{probans}

\begin{probans}
$\alpha = \sqrt{3+a\sqrt{5}}$とする. $K\supset \mb{Q}(\sqrt{5})$は明らかである. まず$K\neq \mb{Q}(\sqrt{5})$を示す. もしそうなら, $\alpha^2 = (p+q\sqrt{5})$となる$p,q\in\mb{Q}$がある. これより$p^2 + 5q^2 = 3$, $pq = a/2$を得る. $p^4 + 5(a/2)^2 = 3p^2$より$(2p)^4 -12(2p)^2 + 20a^2 = 0$だから$\mb{Z}$の整閉性より$2p\in \mb{Z}$. これを$(2p)^2$の二次方程式と見た判別式が$144- 80a^2$なので, $a>1$では判別式が負で不適. $a=1$では判別式は$64$で$(2p)^2 = 10, 2$だが, いずれも不適. よって$[K:\mb{Q}] > 2$で, $\mb{Q}(\sqrt{5})$を中間体に含むので, $[K:\mb{Q}]\geq 4$が分かる. \par
$(\alpha^2 - 3)^2 = 5a^2$で$\alpha$は$f(T) = T^4 - 6T^2 + 9 -5a^2$の根であり, $[K:\mb{Q}]\leq 4$が分かるので, $[K:\mb{Q}]=4$であり, $f$が$\alpha$の最小多項式でなければならない. $f(T)$の$\mb{Q}$上の共役元は$-\alpha, \beta:= \sqrt{3-a\sqrt{5}}, -\beta$であるが, これらが$K$に含まれるためにはすべて実数でなければならないので, $\beta\in \mb{R}$より$3 \geq a\sqrt{5}$でなければならない. よって$a=1$が必要. 逆に$a=1$なら$\alpha \beta = 2$より$\beta = 2/\alpha\in \mb{Q}(\alpha)$なので$K/\mb{Q}$は4次Galois拡大体. よって\bolm{a=1}. \par
ガロア群を$G$とする. $G$は位数4なのでAbel群. $\sigma\in G$は$\sigma(\alpha)\in\set{\alpha,\-\alpha,\beta,-\beta}$で決まり, 位数と組み合わせが4つなのですべての選択が可能. よって$\tau_1(\alpha) = -\alpha$, $\tau_2(\alpha) = \beta$, $\tau_3(\alpha)=-\beta$なる$\tau_i\in G$が存在し, $G=\set{\mr{id},\tau_1,\tau_2,\tau_3}$である. 明らかに$\tau_1^2 = \mr{id}.$ また, $\alpha\beta = 2$なので$\tau_2(\beta) = \alpha$, $\tau_3(-\beta) = \alpha$が分かるから, $\tau_i$はすべて位数2. よって$G\cong \cyc{2}\times \cyc{2}$.
\end{probans}

\begin{probans} %H29専門
(1): $\eta\in N_{S_{5}}(G) \iff \eta G\eta^{-1} = G \iff \lan (\eta(1) \eta(2) \eta(3)) \ran = \lan (123) \ran \iff (\eta(1) \eta(2)\eta(3)) = \sigma, \sigma^2$. よって, $\eta$が$4,5$を固定すれば$\eta\in S_3$で$\eta$が4,5を交換すれば$\eta\tau^{-1} \in S_3$なので, $N_{S_5}(G) = \set{\alpha\tau^k \mid  \alpha\in S_3, k=0,1}$と表せる. $S_3$の元と$\tau\in H$は可換であることに注意. よって$\beta\in S_3$, $\alpha\tau^k\in N_{S_5}(G)$に対し$\alpha\tau^k \beta \tau^{-k} \alpha^{-1} = \alpha\beta\alpha^{-1} \in S_3$だから$S_3$は$N_{S_5}(G)$で正規. $\alpha\tau^k\tau \tau^{-k}\alpha^{-1} = \tau\in H$なので$H$も$N_{S_5}(G)$で正規. よって$S_3H$は群をなし, $S_3\cap H = \set{1}$だから直積に同型. \qed \\
(2): $f(X) = (X-\alpha_1)\dots (X-\alpha_5)$とせよ. $K/\mb{Q}$はGalois拡大. ガロア群を$G$とすると$\sigma\in G$は$\set{\alpha_1,\dots, \alpha_5}$の置換の様子から決定される. とくに単射群準同型$G\hookrightarrow S_5$があるので, $G$を$S_5$の部分群とみなす($\alpha_i$を文字$i$に見立てる). ガロア対応により$F$に対応する部分群を$S$とする. $[K:F] = 3$だから$S$は位数3の部分群. それは位数3の元で生成されるが, $S_5$の元のサイクル分解を考えるとその形は3次巡回置換の形でなければならない. 適当に文字を入れ替えることで, $S = \lan (123) \ran$であると仮定できる. \ao{また, $\sigma\in G$に対して$\sigma(F)$も$[K:\sigma(F)] = 3$を満たすので$F$の一意性から$\sigma(F) = F$である.} ガロア対応において$\sigma(F)$は$\sigma S \sigma^{-1}$に対応するので$\sigma S \sigma^{-1}  =  S$が従う. よって$\sigma \in N_{S_5}(S) = N_{S_5}(\lan (123)\ran) = S_3\times H$ ($H = \lan (45) \ran$)である. これにより$G\subset N_{S_5}(G)$が従う. \par
さて, $f(X)$を割り切る既約三次多項式を見つけるには次のような相異なる$\alpha_i,\alpha_j,\alpha_k$を探すとよい:
\[\alpha_i + \alpha_j + \alpha_k, \alpha_i\alpha_j + \alpha_j \alpha_k + \alpha_k\alpha_i, \alpha_i\alpha_j\alpha_k \in \mb{Q}, \alpha_i,\alpha_j,\alpha_k \notin \mb{Q}\]
実際, これらが満たされれば$(X-\alpha_i)(X-\alpha_j)(X-\alpha_k)$が$\mb{Q}$上既約であるからだ. $s_1 = \alpha_i + \alpha_j + \alpha_k$, $s_2 = \alpha_i\alpha_j + \alpha_j\alpha_k + \alpha_k \alpha_i$, $s_3 = \alpha_i\alpha_j\alpha_k$とおくと, このような$\alpha_i,\alpha_j,\alpha_k$を探すことは次を満たすものを探すことと同じである:
\lefba{
任意の$\sigma\in \gal{(K/\mb{Q})}$に対して$\sigma$は$s_1,s_2,s_3$を固定するが, $\alpha_i,\alpha_j,\alpha_k$はそれぞれあるガロア群の元によって固定されない.
}
そして今の場合, $i=1,j=2,k=3$に取れることを示そう. $s_1,s_2,s_3$が$N_{S_5}(S) = S_{3}\times H$の元で固定されることは明らかなので, $G\subset N_{S_5}(S)$により$G$のすべての元で固定される. $G=\gal{K/\mb{Q}}$は$S$を部分群に持つとしているので, $\sigma := (123)\in G$である. $\alpha_1,\alpha_2,\alpha_3$のどれか一つがすべての$G$の元で固定されているとして矛盾を導く. たとえば$\alpha_1$が固定されるものだとすると, $\sigma,\sigma^2$でも固定されるので
\[\alpha_3 = \sigma^2(\alpha_1) = \alpha_1 = \sigma(\alpha_1) = \alpha_2 \]
なので$\alpha_1=\alpha_2=\alpha_3$である. $s_1 = 3\alpha_1 \in \mb{Q}$より$\alpha_1\in \mb{Q}$なので$f(X) = (X-\alpha_1)^3 (X=\alpha_4)(X=\alpha_5)$なので$K=\mb{Q}(\alpha_4,\alpha_5)$である. $f(x)$から3乗因子を割った $X^2 - (\alpha_4+\alpha_5)X + \alpha_4\alpha_5$も$\mb{Q}$係数であり, とくに $\alpha_4+\alpha_5 \in \mb{Q}$なので$K=\mb{Q}(\alpha_4)$であるが, $\alpha_4$が二次の$\mb{Q}[X]$の元の根であるということは$K/\mb{Q}$が高々2次であるということになるから$[K:F] = 3 > 2$なる$F$の存在に反する. したがって$\alpha_1,\alpha_2,\alpha_3$のいずれも, あるガロア群の元により固定されないことが従う.\qed
\end{probans}



\begin{probans}
(1): $\alpha = \sqrt{1+\sqrt{2}}$とする. $(\alpha^2 - 1)^2 - 2 = 0$より$T^4 - 2T^2 - 1 = 0$. これは$T=S+1$とおくと$2$に関するEis多項式になり既約である(2次と4次だからこの変換で展開したときに偶数係数が現れることと, $T=1$を代入して$-2$になることから見つけられる). よって$\mb{Q}(\alpha)/\mb{Q}$は4次で, $\mb{R}$に含まれているので$K/\mb{Q}$は8次. $\alpha$の$\mb{Q}$上の共役元は, $-\alpha$, $\beta:= \sqrt{\sqrt{2}-1}i , -\beta$である. $\alpha\beta = i$なので$\beta = i/\alpha \in \mb{Q}(\alpha,i)$であり, $K/\mb{Q}$はGalois拡大. ガロア群を$G$とすると$G$は位数8.  $K=\mb{Q}(\alpha,-\alpha,\beta,-\beta)$だから, $\sigma\in G$は$X=\set{\alpha,\beta,-\alpha,-\beta}$を置換し, 逆にその置換で$\sigma$は決定される. これにより$G\subset \maf{S}_{4}$と思える. \aka{$\alpha,\beta,-\alpha,-\beta$をこの順番に$1,2,3,4$とみなす.} $\maf{S}_{4}$の$H=\lan(1234), (24)\ran\cong D_{4}$と$G$が同型であることを示す. $\sigma\in G$は$\alpha,i$での値によっても決定されることに注意する. $\tau(\alpha) = \beta$, $\tau(i) = -i$なる$\tau \in G$を取る. $\alpha\beta = i$なので$\tau(\beta) = -\alpha$であり, $\alpha\xmapsto{\tau} \beta\xmapsto{\tau}-\alpha \xmapsto{\tau} -\beta$だから$\tau = (1234)$である. また, $\eta\in G$を$\eta(\alpha) = \alpha$, $\eta(i) = -i$となるものとすると$\eta(\beta) = -\beta$だから, $\eta$は$\beta$と$-\beta$のみ交換する. よって$\eta=(24)$であるから$H\subset G$とみなせる. $G$は位数8で, $H$の位数も8なので$G=H$である. \qed \\
(2): $K,\mb{Q}$は自明な中間体. Galois対応により$G\cong D_{4}$の部分群を見る. 指数4のものは位数2であり位数2の元によって生成される. $H=\set{1, (1234),(13)(24),(1432),(24), (12)(34), (13), (23)(14), }$であり, 位数2の元は4-cycleと単位元以外で5個ある.\par  $\sqrt{2} = \alpha^2 - 1$に注意すると,   次のように4次拡大体が得られる.
\bealn{
\lan (13)(24)\ran  の固定体は \mb{Q}(\sqrt{2}, i)\\
\lan (24) \ran の固定体は \mb{Q}(\alpha) \\
\lan (12)(34) \ran の固定体は \mb{Q}(\alpha+\beta)\\
\lan (13) \ran の固定体は\mb{Q}(\beta) \\
\lan (23)(14) \ran の固定体は\mb{Q}(\alpha-\beta)
}
最後に指数2の中間体を見つける. $\eta=(1234)$は, $\eta(i) = -i$と$\eta(\sqrt{2}) = \beta^2 - 1 = -\sqrt{2}$を満たすので$\sqrt{-2}$を固定する. よって
\bealn{
\lan (1234) \ran = \lan (1432) \ran の固定体は = \mb{Q}(\sqrt{-2})
}
$\mb{Q}(i),\mb{Q}(\sqrt{2})$という中間体から部分群を見つける. $\sqrt{2} = \alpha^2 - 1$だから$\alpha$を$\set{\alpha,-\alpha}$に移すものが$\mb{Q}(\sqrt{2})$を固定する. それらは$(13)(24), (13), (24), 1$だから
\[\set{1,(13),(24),(13)(24))} の固定体は\mb{Q}(\sqrt{2})\]
$i=\alpha\beta$なので, $\set{\alpha,\beta}$が$\set{\alpha,\beta}$に移るものや$\set{-\alpha,-\beta}$に移るものは$\mb{Q}(i)$を固定する. そのようなものは$1,(13)(24),(14)(23),(12)(34)$であり, したがって
\bealn{
\set{1,(13)(24),(14)(23),(12)(34)}の固定体は\mb{Q}(\sqrt{2}).
}
最後に, $G$の位数4の部分群$K$がこれらに限ることを示す. $K$が巡回群なら明らかに$K\ni \lan (1234)\ran$である. $K$が巡回群でないとすれば, $K$の単位元でない元の位数は必ず2である. $K\subset \set{1,(13)(24),(12)(34), (23)(14), (13),(24)}$である. $(13)(24)\in K$の場合, $(13)\in K \iff (24)\in K$であり, $(12)(34)\in K\iff (23)(14)\in K$が分かるので, すでに挙がっている群が得られる. $(13)(24)\notin K$とせよ. このとき$K$は必ず互換$\beta$を含み, 位数2の偶置換$\alpha$も含む. このとき$\alpha\beta$は$\beta$と異なる奇置換である. $\lan (1234) \ran$は互換を含まないので$\alpha\beta$が互換であることが分かる. よって$K$は$(13),(24)$を含まねばならず, $(13)(24)\in K$なので仮定に反する. よって位数4の部分群はこれらで尽くされる.
\end{probans}

\begin{probans} %R2senmon
$\alpha = \sqrt{-p} + \sqrt[3]{p}$とする. \[(\alpha - \sqrt{-p})^3 = \alpha^3 -3p\alpha - \sqrt{-p}(3\alpha^2 -p)  = p\]
より$\sqrt{-p}$が$\alpha$と$p$の有理式で書ける. すなわち$\sqrt{-p} \in \mb{Q}(\alpha)$である. よって$\sqrt[3]{p}\in \mb{Q}(\alpha)$も従う. $[\mb{Q}(\sqrt[3]{p}, \sqrt{p}):\mb{Q}] = 6$なので$\alpha$の最小多項式は6次. それは先の式で平方を取ったものから構成できる.すなわち$\alpha$の最小多項式は
\[(T^3 - 3pT - p)^2 +p(3T^2 - p)^2 = 0\]
である. $\alpha_1 = \sqrt{-p}+\zeta_3\sqrt[3]{p}$, $\alpha_2 = \sqrt{-p} + \zeta_3^2\sqrt[3]{p}$もこの方程式を満たすので, これは$\alpha$の$\mb{Q}$上共役元. よって$K\ni \alpha_1,\alpha_2$である. また,
\[\alpha_2 - \alpha_1 = \zeta_3(\zeta_3-1)\sqrt[3]{p}, \alpha_1-\alpha = (\zeta_3 - 1)\sqrt[3]{p}\]
も$K$に含まれるので, $\frac{\alpha_2-\alpha_1}{\alpha_1-\alpha} = \zeta_3\in K$. また,
\[\alpha + \alpha_1 + \alpha_2 = 3\sqrt{-p}\in K\implies \sqrt{-p}\in K\]
であり, $\alpha - \sqrt{-p} = \sqrt[3]{p}\in K$も得られる. よって, $K\ni \sqrt{-p}, \sqrt[3]{p}, \zeta_3$である. \par
ここで, $\beta_{t} = - \sqrt{-p} + \zeta_3^{t}\sqrt[3]{p}$とおくと, $\beta_{t}$も先と同じ6次方程式を満たす. なぜなら, たとえば $\beta=\beta_0$の場合は
\[(\beta + \sqrt{-p})^3 = \beta^3 - 3p\beta + \sqrt{-p}(3\beta^2 - p) = p\]
となり, $\sqrt{-p}$の前で符号が(先の$\alpha$のものと)異なるだけで, 2乗すれば同じ式が得られるから. よって, $\alpha,\alpha_1,\alpha_2,\beta,\beta_1,\beta_2\in \mb{Q}(\sqrt{-p}, \sqrt[3]{p},\zeta_3)$である. この体は明らかにGalois核大体である($\alpha$の共役元はこの6個の中にある). よって $K$はこの体に含まれる. 逆に任意の$\alpha$を含むGalois拡大体は共役元を含み, この体を含まざるを得ないので$K$はこの体を含む. よって$K=\mb{Q}(\sqrt{-p},\sqrt[3]{p},\zeta_3)$を得る. \par
さて, $K$を二つのGalois拡大体の合成体$\mb{Q}(\sqrt{-p})\cdot \mb{Q}(\sqrt[3]{p},\zeta_3)$とみる. まず, 6次核大体$\mb{Q}(\sqrt[3]{p},\zeta_3)$のガロア群$G_2$について論じる. $\sigma\in G_2$なら$\sigma(\sqrt[3]{p})\in \set{\sqrt[3]{p},\sqrt[3]{p}\zeta_3,\sqrt[3]{p}\zeta_3^2}, \sigma(\zeta_3)\in \set{\zeta_3,\zeta_3^2}$で決定され, 6次なのでこの選択はすべて実現できる. $\sigma(\sqrt[3]{p})$と$\sigma(\sqrt[3]{p}\zeta_3)$から$\sigma(\zeta_3)$も決まるので, $\sigma$は$\set{\sqrt[3]{p},\sqrt[3]{p}\zeta_3,\sqrt[3]{p}\zeta_3^2}$をどのように置換するかという情報からも決定される. $G_2$がこの3点集合$S$に作用すると考え, $G_2\to \mr{Aut}(S)\cong \maf{S}_{3}$を得る. これが位数6の群の間の単射準同型であるため, 同型である. \\
次に$K_0 = \mb{Q}(\sqrt{-p})\cap \mb{Q}(\sqrt[3]{p},\zeta_3)$とする. $[K_0:\mb{Q}]\leq 2$は明らかである.
\lefba{ 仮にこれが2であったとせよ. $\mb{Q}(\sqrt[3]{p},\zeta_3)$の中間体$K_0$は$\maf{S}_{3}$の指数2の部分群に対応する. $\maf{S}_{3}$の指数2の部分群は一つであるから, $K_0 = \mb{Q}(\zeta_3)$が従う. よって$\zeta_3\in \mb{Q}(\sqrt{-p})$なので, $\sqrt{-p} = a\zeta_3 + b$と置いて, $0 = a\zeta_3^2 + 2ab\zeta_3 + b^2 + p$だから, $\zeta_3^2 + \zeta_3 + 1 = 0$と合わせると$a=2ab = b^2+p$. よって$a=0$だが$b^2+p=0$ではなく矛盾. (あるいはKummer理論を使ってもよい)}
よって$[K_0:\mb{Q}] = 1$より$K_0=\mb{Q}$を得る. Galois推進定理より
\[\mr{Gal}(K/\mb{Q})\cong \mr{Gal}(\mb{Q}(\sqrt{-p})/\mb{Q}) \times \mr{Gal}(\mb{Q}(\sqrt[3]{p},\zeta_3)/\mb{Q}) \cong \cyc{2}\times \maf{S}_{3}\]


\end{probans}




}
\newpage




\section{補足}
\subsection{アルティン・シュライアー拡大}
\begin{egprob} (2020代数学II, 12/8, 練習問題)
$p$を素数とする. $L/K$を標数$p$の体の代数拡大とする. $\alpha\in K$, $\beta\in L$は$\beta^p - \beta - \alpha = 0$を満たすと仮定する.
\ben
\item $T^p - T - \alpha$を$K(\beta)[T]$において因数分解せよ. (\tf{Hint:} $T=\beta + i$ ($0\leq i\leq p-1$)を代入せよ. )
\item $\beta\notin K$とする. このとき$K(\beta)/K$は$p$次ガロア拡大であって, そのガロア群は$\mb{Z}/p\mb{Z}$と同型であることを示せ.
\een
\end{egprob}









